%

%\documentclass[11pt,a4paper,twoside]{article}

\usepackage[utf8]{inputenc}

\usepackage[T1,T2A]{fontenc}

\usepackage[english.ancient,russian.ancient]{babel}

\usepackage{graphicx}

\usepackage{array}

\usepackage[doi]{samgtu-bib}

\usepackage{samgtu}

\usepackage{samgtu-name}

\usepackage{mathtools}

\mathtoolsset{showonlyrefs}

\usepackage{desclist}

\usepackage{euscript}

\usepackage{mathrsfs} 

\usepackage{multicol}

\usepackage{mathrsfs}

\usepackage{cite}

\usepackage{cmap}

\usepackage{qrcode}

\allowdisplaybreaks[4]

\usepackage[nodayofweek]{datetime}

\usepackage[thinlines]{easybmat}



\renewcommand{\JournalYear}{2018}

\renewcommand{\JournalVolume}{22}

\renewcommand{\JournalNumber}{4}



\newdate{JournalDateSign}{29}{12}{2018}% Дата подписания Вестника в печать

\renewcommand{\JournalDateSign}{\displaydate{JournalDateSign}}

\newdate{JournalDatePrint}{16}{01}{2019}% Дата выхода Вестника из печати

\renewcommand{\JournalDatePrint}{\displaydate{JournalDatePrint}}

%\renewcommand{\JournalUPL}{16.56} % Усл. печ. л.

%\renewcommand{\JournalUIL}{16.53} % Уч.-изд. л.

\renewcommand{\JournalUPL}{15.24} % Усл. печ. л.

\renewcommand{\JournalUIL}{15.21} % Уч.-изд. л.

\renewcommand{\JournalRegNumber}{220/18} % Регистрационный номер

\renewcommand{\JournalOrderNumber}{2} % Номер заказа



%\renewcommand{\JournalRMonth}{Март} % Месяц выхода Вестника

%\renewcommand{\JournalEMonth}{March} % Месяц выхода Вестника

%\renewcommand{\JournalRMonth}{Июнь} % Месяц выхода Вестника

%\renewcommand{\JournalEMonth}{June} % Месяц выхода Вестника

%\renewcommand{\JournalRMonth}{Сентябрь} % Месяц выхода Вестника

%\renewcommand{\JournalEMonth}{September} % Месяц выхода Вестника

\renewcommand{\JournalRMonth}{Декабрь} % Месяц выхода Вестника

\renewcommand{\JournalEMonth}{December} % Месяц выхода Вестника



\newcommand{\totop}[1]{\raisebox{6pt}[1pt][2pt]{#1}}

\newcommand{\tottop}[1]{\raisebox{12pt}[1pt][2pt]{#1}} 

\newcommand{\totttop}[1]{\raisebox{18pt}[1pt][2pt]{#1}} 

\newcommand{\tottttop}[1]{\raisebox{24pt}[1pt][2pt]{#1}} 

\newcommand{\totttttop}[1]{\raisebox{30pt}[1pt][2pt]{#1}} 

\newcommand{\tobot}[1]{\raisebox{-6pt}[1pt][2pt]{#1}} 

\newcommand{\tobbot}[1]{\raisebox{-12pt}[1pt][2pt]{#1}} 

\newcommand{\tobbbot}[1]{\raisebox{-18pt}[1pt][2pt]{#1}}



\graphicspath{{./00PIC/}}

%

%\usepackage{samgtu-name}

%

%\usepackage{qrcode}

%

%%

%%\newcommand{\totop}[1]{\raisebox{6pt}[1pt][2pt]{#1}}

%%\newcommand{\tottop}[1]{\raisebox{12pt}[1pt][2pt]{#1}}

%%\newcommand{\totttop}[1]{\raisebox{18pt}[1pt][2pt]{#1}}

%%\newcommand{\tottttop}[1]{\raisebox{24pt}[1pt][2pt]{#1}}

%%\newcommand{\totttttop}[1]{\raisebox{30pt}[1pt][2pt]{#1}}

%\newcommand{\tobot}[1]{\raisebox{-6pt}[1pt][2pt]{#1}}

%%\newcommand{\tobbot}[1]{\raisebox{-12pt}[1pt][2pt]{#1}}

%%\newcommand{\tobbbot}[1]{\raisebox{-18pt}[1pt][2pt]{#1}}

%

%%\graphicspath{{./00PIC/}}

%

%%\samgtupubl % для генерации pdf, отдаваемого в типографию СамГТУ

%\pdfcrop % обрезка pdf для размещения на math-net.ru

%\draft % На корректуре

%%\briefcomm % Статья является кратким сообщением

%

\vsgtu{1566}

%

\FirstPage{116}

\LastPage{152}

%

%

\ResearchArticle

%

%%

%%

%\begin{document}







\renewcommand{\baselinestretch}{0.9}



\hfill \qrcode[hyperlink,height=.8in]{http://doi.org/10.14498/vsgtu1566}

\vspace{-.8in}





% \Partition{Механика деформируемого твердого тела}{Applied mechanics} % Раздел Вестника, в который подается статья

% Названия разделов см. на сайте при подаче заявки





\udc{536.2:46(075.8)}

\msc{35K05, 80A20}



\rutitle{Разработка математических моделей и~исследование неравновесных явлений с~учетом пространственно-временной нелокальности}

\rutitleshort{Разработка математических моделей и исследование неравновесных явлений\dots}

\rutitlehack{Разработка математических моделей \\ и~исследование~неравновесных явлений \\ с~учетом~пространственно-временной \\ нелокальности \vspace{-.8mm}}



\ruauthor{И.~В.~Кудинов}{К\,у\,д\,и\,н\,о\,в~И.\,В.}



\ruaffil{\rusamgtu.}

\enaffil{\ensamgtu.}







\ruauthorsdetails{1}{

\fullauthor{\rupid{44183}{Игорь Васильевич Кудинов}}

\CAuthor

\orcid{0000-0002-9422-0367}

\regalia{кандидат технических наук} 

\position{доцент} 

\dept{каф. теоретических основ теплотехники и гидромеханики} 

\email{igor-kudinov@bk.ru}

}



\enauthorsdetails{1}{

\fullauthor{\enpid{44183}{Igor V. Kudinov}} 

\CAuthor

\orcid{0000-0002-9422-0367}

\regalia{Cand. Techn. Sci.}

\position{Associate Professor} 

\dept{Dept. of Theoretical Fundamentals of Heat-Engineering and Hydromechanics}

\email{igor-kudinov@bk.ru}

}







\rukeywords{локально"=неравновесные процессы, пространственно"=временная нелокальность,

коэффициенты релаксации, колебания уп\-ру\-гих тел, жидкостей и газов,

аномальная теплопроводность, динамические температурные напряжения, тепловое

воспламенение}



\ruabstract{На основе принципов локально-неравновесной термодинамики разработаны

математические модели процессов переноса тепла, массы, импульса с учетом

пространственно-временной нелокальности. Вывод дифференциальных уравнений

переноса основывается на учете в диффузионных законах Фурье, Фика, Ньютона,

Гука, Ома ускорения во времени как удельных потоков (тепла, массы,

импульса), так и градиентов соответствующих величин. Исследования точных

аналитических решений полученных моделей позволили обнаружить новые

закономерности изменения искомых параметров для малых и сверхмалых значений

временной и пространственной переменных, а также для быстропротекающих

процессов, время изменения которых сопоставимо со временем релаксации. И, в

частности, из анализа точного аналитического решения обнаружен факт задержки

во времени принятия граничного условия первого рода, свидетельствующий о

том, что ввиду сопротивления тела, оказываемого процессу проникновения

теплоты, его мгновенный прогрев на границе невозможен ни при каких условиях

теплообмена с окружающей средой. Следовательно, коэффициент теплоотдачи на

стенке зависит не только от условий теплообмена (скорость среды, вязкость и~прочее), но и от физических свойств тела, и он, во-первых, является переменным

во времени и, во-вторых, не может превысить некоторой

предельной для каждого конкретного случая величины.}



\entitle{Development of mathematical models and research of strongly nonequilibrium phenomena taking into account space-time nonlocality}

\entitleshort{Development of mathematical models and research of strongly nonequilibrium phenomena\dots}



\enauthor{I.~V.~Kudinov}{K\,u\,d\,i\,n\,o\,v~I.~V.}





\enabstract{Based on the principles of locally-nonequilibrium thermodynamics, the mathematical models of heat, mass, momentum transfer processes were developed   taking into account spatial and temporal nonlocality. The derivation of the differential equations is based on the taking into account accelerate in time as the specific fluxes (heat, mass, momentum) and the gradients of the corresponding variables in the Fourier's, Fick's, Newton's, Hooke's, Ohm's, etc. diffusion laws. Studying of exact analytical solutions of the obtained models allowed us to discover new regularities of the changes of the desired parameters at low and ultra low values of temporal and spatial variables, and for all fast processes, time change which is comparable with the relaxation time. And, in particular, from the analysis of the exact analytical decision the fact of a time lag of acceptance of a boundary condition of the first kind demonstrating that in view of resistance of the body shown to warmth penetration process, its instantaneous warming up on boundary is impossible under no circumstances heat exchange with the environment is found. Therefore, the heat emission coefficient on a wall depends not only on heat exchange conditions (environment speed, viscosity and so forth), but also on physical properties of a body and it, in the first, is variable value in time and, in the second, it can not exceed some value, limit for each case.}



\enkeywords{local nonequilibrium heat transfer processes, spatial and temporal

nonlocality, relaxation coefficients, vibrations of elastic bodies, fluids

and gases, abnormal conduction, dynamic thermal stresses, thermal ignition}



\newdate{kuda}{17}{10}{2017}

\date{\displaydate{kuda}}

\newdate{kudb}{10}{02}{2018}

\revisiondate{\displaydate{kudb}}

\newdate{kudc}{12}{03}{2018}

\accepteddate{\displaydate{kudc}}



\newdate{kude}{31}{03}{2018}

\renewcommand{\JournalDataOnline}{\displaydate{kude}}





\makerutitle



\smallskip





\Section[n]{Введение}

Изучение различного рода локально"=неравновесных систем \linebreak и~процессов переноса

(энергии, массы, импульса или их аналогов) основано на различных моделях:

локально"=неравновесной термодинамики 

(Sybren R.~de Groot, Peter Mazur~\cite{kud:bib1},

Istv\'an Gyarmati~\cite{kud:bib2},

William~A. Day~\cite{kud:W.Day}, Nikola Petrov, Jordan G.~Brankov~\cite{kud:Brankov},

David~Jou,  Jos\'e Casas-V\'azquez, Georgy Lebon~\mbox{\cite{kud:Lebon, kud:bib5}}, 

С.~Л.~Соболев~[\citen{kud:bib3}--\citen{kud:bib43}]); 

кинетических (James C.~Maxwell~\cite{kud:Maxwell}, 

В.~А.~Бубнов~\cite{kud:bib41}, 

Borivoje~B. Mikic~\cite{kud:Mikic},

D.~Urushev, M.~Borisov, A.~Vavrek~\cite{kud:Urushev}); 

феноменологических 

(Morton~E. Gurtin, A.~C.~Pipkin~\cite{kud:Gurtin},

Jace~W. Nunziato~\cite{kud:Nunziato}, 

Helmut~D. Weymann~\cite{kud:Weymann}, 

Yehuda Taitel~\cite{kud:Taitel}, 

Е.~И.~Леванов, Е.~Н.~Сотский~\cite{kud:bib31}). 

В~некоторых из перечисленных направлений рассматривается понятие <<тепловой

памяти>>, среды с внутренними степенями свободы, среды скоростного типа.

Существуют методы, основанные на предположении, что удельная энтропия

является одновременно функцией равновесных экстенсивных параметров и их

потоков "--- <<скоростных переменных>>, описывающих инерционные свойства

систем при установлении локального равновесия. Уравнения

локально"=неравновесного переноса получают также из уравнения Больцмана путем

использования теории случайных блужданий, а также молекулярно"=кинетическими

методами.



В различных версиях перечисленных теорий решается вопрос о том, что понимать

под температурой локально"=неравновесных систем. В классической термодинамике

температура определяется только для равновесных или локально"=равновесных

систем. В различных версиях локально"=неравновесной термодинамики вводятся

понятия локально"=неравновесных энтропии и внутренней энергии, и

локально"=неравновесная температура определяется как производная от

внутренней энергии по энтропии (David~Jou,  Jos\'e Casas-V\'azquez, Georgy Lebon~\mbox{\cite{kud:Lebon, kud:bib5}}). В~средах скоростного типа, когда состояние системы зависит от скорости

изменения температуры, абсолютная температура заменяется на

<<термодинамическую>> (Nikola Petrov, Jordan G.~Brankov~\cite{kud:Brankov}) 

Обе эти температуры совпадают,

если абсолютная температура неизменна во времени. В случае, когда

температура изменяется во времени, термодинамическая температура

представляется в виде разложения в ряд Тейлора по степеням абсолютной

температуры.



Так как величина абсолютной температуры пропорциональна скорости ее

изменения, величина отклонения термодинамической температуры от

абсолютной характеризует степень локальной неравновесности системы. 

В~молекулярной динамике вводится понятие <<кинетической>> температуры 

как меры средней кинетической энергии молекул, определяемой по

локально"=неравновесной функции распределения 

(С.~Л.~Черешнев, А.~П.~Генич, С.~В.~Куликов, Г.~Б.~Манелис~\cite{kud:bib32},

Rosemary~A. MacDonald,  Donald~H. Tsai~\cite{kud:MacDonald}).



Необходимо отметить, что из достаточно большого числа версий

локально"=неравновесных процессов некоторые из них приводят к противоречивым

результатам. Так, устраняя проблемы параболических уравнений, связанные с~бесконечной скоростью распространения потенциалов исследуемых полей,

появляются новые проблемы, связанные со скачками искомой функции, с~отрицательными ее значениями в~обратной волне, а также с бесконечными

величинами искомой функции в обратной волне.



В настоящее время наиболее перспективным, по мнению автора данной статьи,

является направление исследований, связанное с использованием не только

временного, но и пространственного микромасштаба релаксации. Для его оценки

в некоторых теориях вводится <<поток теплового потока>>, характеризующий

влияние пространственного микромасштаба системы путем учета длины свободного

пробега микрочастиц.



Несмотря на наличие большого числа версий локально"=неравновесных процессов

переноса тепла, массы, импульса, в настоящее время отсутствует их единая

теоретическая база, на основе которой можно было бы построить

непротиворечивые теории по любому из направлений переноса.



В последнее время в связи с необходимостью проведения различных физических

процессов в экстремальных условиях возник интерес к изучению

локально"=неравновесных процессов. Интерес к таким процессам обусловлен

возможностями познания новых закономерностей и свойств материалов, которые

могут быть широко использованы для создания новых технологий получения

наноматериалов и покрытий с уникальными физико"=химическими свойствами.

Однако для описания процессов, протекающих в локально"=неравновесных

условиях, классическая термодинамика неравновесных процессов не может быть

использована, так как она основана на принципе локального равновесия,

справедливого для относительно медленных процессов. Возникает проблема

создания адекватных физических и соответствующих математических моделей

описания локально"=неравновесных процессов на основе современной версии

термодинамики третьего типа "--- расширенной необратимой термодинамики,

выходящей за рамки принципа локального равновесия. %[\citen{kud:bib1}--\citen{kud:bib14}, \citen{kud:bib24}--\citen{kud:bib44}].



\smallskip



\Section{Основные положения предлагаемого метода исследования}

Известные математические постановки краевых задач математической физики основаны на бесконечной скорости распространения потенциалов исследуемых полей (задачи теплопроводности, теплообмена в жидкостях, динамической

термоупругости, колебаний упругих тел и жидкостей и~проч.), которая заложена в соответствующих диффузионных уравнениях: закон Фурье (для теплового потока), Фика (для задач диффузии), Гука (для нормального

напряжения), Ньютона (для касательного напряжения в жидкости), Ома (для распространения тока в проводниках) и~др. Полученные на основе этих законов дифференциальные уравнения являются локальными, то есть в них не учитывается

пространственно"=временная нелокальность. 

Например, математические постановки классических задач нестационарного тепломассопереноса основаны на

параболических уравнениях, вывод которых базируется на уравнении теплового баланса (законе сохранения энергии) и феноменологическом законе Фурье для теплового потока (гипотезе Фурье):

\begin{equation}

\label{kud:eq1}

q=-\lambda \frac{ \partial T}{\partial x}.

\end{equation}

Здесь $q$ "--- тепловой поток; $\lambda$ "--- коэффициент теплопроводности; $T$

"--- температура; $x$ "--- координата.



Исследования аналитических решений полученных таким путем уравнений приводят

к заключению о том, что скорость распространения тепла в~изучаемых средах

является бесконечной. Для большинства практических задач нестационарной

теплопроводности классическая феноменология Фурье приводит к

удовлетворительным результатам, подтверждаемым опытными данными. Однако

существует большой класс задач, для которых классическая модель

теплопроводности оказывается несостоятельной. Это, например, краевые задачи

при малых и сверхмалых значениях временной и пространственной переменных,

задачи для быстропротекающих процессов, время протекания которых сопоставимо

со временем релаксации, и др. 

Основная причина неприменимости классической

модели к указанным процессам в том, что закон Фурье является лишь

приближенным описанием процесса теплопроводности, не учитывающим

инерционность теплового потока и скалярной величины градиента температуры.

Первыми, кто ввел инерционность в уравнения переноса, были Максвелл,

Вернотт, Каттанео. 

В~частности, ими была предложена формула для теплового потока, известная как формула Максвелла"--~Каттанео, путем использования в~законе Фурье релаксационного слагаемого, учитывающего конечную скорость распространения теплоты [\citen{kud:bib5}--\citen{kud:bib44},  \citen{kud:bib6}--\citen{kud:bib8}]:

\begin{equation}

\label{kud:eq2}

q=-\lambda \frac{\partial T}{\partial x}-\tau _1 \frac{\partial q}{\partial

t}.

\end{equation}

Здесь $t$ "--- время; $\tau _1$ "--- коэффициент релаксации теплового потока $q$.



Подставлением \eqref{kud:eq2} в уравнение теплового баланса

\begin{equation}

\label{kud:eq3}

c\rho \frac{ \partial T}{\partial t}= \frac{\partial q} {\partial x},

\end{equation}

где

$c$ "--- теплоемкость; $\rho $ "--- плотность, выводится классическое гиперболическое уравнение теплопроводности в безразмерном виде: %(в~электротехнике оно известно как телеграфное), содержащее, в~отличие от параболического, вторую производную по времени от искомой функции:

\begin{equation}

\label{kud:eq4}

\frac{\partial \Theta(\xi, \operatorname{\sf Fo})}{\partial

\operatorname{\sf Fo}}+\operatorname{\sf Fo}_1 \frac{\partial^2\Theta (\xi , \operatorname{\sf Fo})}{\partial \operatorname{\sf Fo}^2}=\frac{\partial ^2 \Theta

(\xi, \operatorname{\sf Fo})}{\partial \xi^2}.

\end{equation}

Здесь

$$ \Theta=\frac{T-T_{w}}{T_0 -T_{w}},\quad 

\operatorname{\sf Fo}=\frac {at}{ \delta^2} ,\quad 

\xi=\frac x \delta;$$

%{\color{red}  

$T$ "--- температура, $T_{w}$ "--- температура стенки, $T_{0}$ "--- начальная температура, $t$ "--- время,

$a$ "--- коэффициент температуропроводности, 

$\delta $ "--- определяющий размер, например, толщина пластины, если уравнение \eqref{kud:eq4} рассматривать как уравнение теплопроводности для пластины; $\Theta$ "--- безразмерная температура, $\operatorname{\sf Fo}$ "--- безразмерное время, $\xi$ "--- безразмерная координата. 

%}









Точное аналитическое решение уравнения \eqref{kud:eq4} при однородных граничных условиях первого рода 

$$

\Theta (0, \operatorname{\sf Fo})=0, \quad \Theta (1, \operatorname{\sf Fo})= 0

$$

и 

начальном условии

$$

\Theta (\xi, 0)=1

$$

имеет вид 

\begin{multline}

\label{kud:eq5}

\Theta (\xi, \operatorname{\sf Fo})=

\sum\limits_{k=1}^\infty 

\Bigl[ 

\exp (s \operatorname{\sf Fo}) \bigl[

B_{1k}  \cos(\beta_k \operatorname{\sf Fo})- \\[-3mm] -

B_{2k}  \sin( \beta_k \operatorname{\sf Fo})

\bigr] \Bigr]  

\cos \Bigl( \frac{2k-1 }{2} (1- 2\xi)\pi \Bigr),

\end{multline}

где 

$$

B_{1k} = \frac{\beta_k B_{2k}} s, \quad 

B_{2k} = \frac{4s } {(2k-1) \beta _k \pi } ,$$

$$

\beta_k = \frac{\sqrt {4 \operatorname{\sf Fo}_1  {\nu}_k -1}} {2\operatorname{\sf Fo}_1}, \quad

s=-\frac1 {4\operatorname{\sf Fo}_1},

\quad \operatorname{\sf Fo}_1 = \frac{a \tau_1}{\delta^2};

$$

%{\color{red}

$\tau_1$ "--- коэффициент релаксации теплового потока;

$\operatorname{\sf Fo}_1 $ "--- безразмерный коэффициент релаксации теплового потока.

%}



Многочисленные исследования решения \eqref{kud:eq5} позволили заключить, что уравнение

\eqref{kud:eq4}, так же как и параболическое уравнение (то есть при $\operatorname{\sf Fo}_{1}

=0$), не обеспечивает достаточно адекватного описания исследуемых процессов.

И,~в~частности, решения гиперболического уравнения \eqref{kud:eq4} приводят к

возникновению скачков температуры, что равнозначно возникновению изотерм

внутри тела и бесконечных тепловых потоков на фронте тепловой волны, а~также

к появлению отрицательных температур в обратной волне. Эти результаты

свидетельствуют о нарушении принципа максимального значения в~краевых

задачах и закона сохранении энергии (рис.~\ref{kud:pic1}). 

%{\color{red} При каких граничных (начальных) условиях решалось уравнение \eqref{kud:eq5}, чтобы получились такие картинки?} 

Анализ причин получаемых из

решения \eqref{kud:eq5} результатов позволил сделать заключение о необходимости учета

инерционности не только теплового потока (временная нелокальность), но и

скалярной величины градиента температуры (пространственная нелокальность).



\begin{figure}[b!]\centering 

\footnotesize

\includegraphics[width=\textwidth]{kudinov1}\\

\caption{Распределение температуры $\Theta$ для различных чисел $\sf Fo$ (значения чисел "--- метки у кривых) в~пластине при $n=10\,000$, где $n$ "--- число членов ряда точного аналитического решения,

для $\operatorname{\sf Fo}_1=0.0015$ (слева) и $\operatorname{\sf Fo}_1=0.1$ (справа)\label{kud:pic1}}

\smallskip



[Figure~\ref{kud:pic1}. The temperature distribution in the plate  for different Fourier numbers ($\sf Fo$)  (the~values of the Fourier numbers are the labels of the curves) when $n=10\,000$ ($n$ is the number~of terms in the series of~the exact analytical solution) for $\operatorname{\sf Fo}_1=0.0015$ (left) and~$\operatorname{\sf Fo}_1=0.1$ (right)]



\end{figure}



Полученная таким путем формула \eqref{kud:eq1} закона Фурье для теплового потока

принимает вид

\begin{equation}

\label{kud:eq6}

q+ \tau_1 \frac{\partial q}{\partial t}+ \tau_1^2

\frac{\partial^2q}{\partial t^2}+\cdots =- 

\lambda

\Bigl(

\frac{\partial T}{\partial x}+ \tau_2 \frac{\partial^2T}{\partial

x\partial t}+ \tau_2^2 \frac{\partial^3T}{\partial

x\partial t^2}+\cdots

\Bigr),

\end{equation}

где $\tau_1 $ "--- коэффициент релаксации теплового потока;

$\tau_2$ "--- коэффициент релаксации скалярной величины градиента

температуры.

Аналогичные уравнения для законов Ньютона $\tau=\mu \partial

\vartheta /\partial y$ и Гука $\sigma=E\partial U/\partial x$ имеют

вид

\begin{gather}

\label{kud:eq8}

\tau + \tau_{1N} \frac{\partial  \tau}{\partial t}+ 

\tau_{1N}^2 \frac{\partial^2 \tau}{\partial t^2}+\cdots=

\mu \Bigl( \frac{\partial \vartheta }{\partial y}+ \tau_{2N}

\frac{\partial ^2\vartheta }{\partial y\partial t}+ \tau_{2N}^2 

\frac{\partial^3 \vartheta }{\partial y\partial t^2}+\cdots \Bigr);\\

\label{kud:eq8}

\sigma+ \tau_{1G} \frac{\partial \sigma }{\partial t}+ 

\tau_{1G}^2 \frac{\partial ^2 \sigma}{\partial t^2}+\cdots=

E\Bigl(\frac{\partial U}{\partial x}+ \tau_{2G} 

\frac{\partial ^2 U}{\partial x\partial t}+ 

\tau_{2G}^2 \frac{\partial^3 U}{\partial x\partial t^2}+\cdots \Bigr),

\end{gather}

где $\tau$, $\sigma$ "--- касательное и нормальное напряжения;

$\vartheta $ "--- скорость; $U$ "--- перемещение; 

$ \mu $ "--- динамическая вязкость;

$E$ "--- модуль упругости; $\tau_{1N}$, $\tau_{1G}$ "---

коэффициенты релаксации напряжений $\tau$ и $\sigma $;  $\tau_{2N}$, $\tau_{2G}$  "--- коэффициенты релаксации  градиентов скорости и деформации $\frac{\partial \vartheta}{\partial y}$ и $\frac{\partial U}{\partial x}$. 



Выражая $q$, $\tau$, $\sigma$, ограничиваясь двумя членами

правой и левой части соотношений \eqref{kud:eq6}--\eqref{kud:eq8}, находим

\begin{gather}

\label{kud:eq9}

q=-\lambda \frac{\partial T}{\partial x}-\lambda \tau

_2 \frac{\partial ^2T}{\partial x\partial t}- \tau_1

\frac{\partial q}{\partial t};\\

\label{kud:eq10}

 \tau= \mu \frac{\partial \vartheta }{\partial y}+ \mu

\tau_{2N} \frac{\partial^2 \vartheta }{\partial y\partial t}- \tau_{1N} \frac{\partial \tau}{\partial t};\\

\label{kud:eq11}

\sigma =E\frac{\partial U}{\partial x}+E \tau_{2G} \frac{\partial ^2U}{\partial x\partial t}- \tau_{1G} \frac{\partial  \sigma}{\partial t}.

\end{gather}



Отметим, что соотношения \eqref{kud:eq9}--\eqref{kud:eq11} совпадают с формулами, полученными из

сформулированной А.~В.~Лыковым системы дифференциальных уравнений Онзагера (Lars Onsager)

(исходя из гипотезы о конечной скорости~диффузии~массы~и~теплоты)~\cite{kud:bib8}:

\begin{equation}

\label{kud:eq12}

J_i =L_i^{(r)} \frac{\partial J_i }{\partial t}+

\sum\limits_{k=1}^N \Bigl(

L_{ik} X_k +L'_{ik} \frac{\partial X_k }{\partial t} \Bigr),

\end{equation}

где $J_i$ "--- поток субстанции (теплоты, массы и т.д.); 

$X_k$ "--- термодинамические движущие силы (градиенты соответствующих величин);

$L_i^{(r)}$, $L_{ik}$, $L'_{ik} $ "--- постоянные феноменологические

коэффициенты переноса $( L_{ik} =L_{ki} )$.



Если положить $L=- \lambda$,

$L^{(r)} =- \tau_1$,

$L'=- \lambda \tau_2$, 

$X_k =\frac{\partial T}{\partial x}$, то из \eqref{kud:eq12}

получаем формулу \eqref{kud:eq9}. 



Если положить $L= {\mu }$,

$L'= \mu \tau _{2N} $, 

$L^{(r)}=- \tau_{1N} $,

$X_k =\frac{\partial \vartheta}{\partial x}$, то

уравнение \eqref{kud:eq12} оказывается тождественным уравнению \eqref{kud:eq10}~\cite{kud:bib8}. 

При $L=E$,

$L'=E\tau _{2G} $, 

$L^{(r)}=-\tau_{1G}$ из \eqref{kud:eq12} получаем формулу \eqref{kud:eq11}.



Соотношения \eqref{kud:eq10}, \eqref{kud:eq11} полностью совпадают с усложненными моделями

вязкоупругого тела, известными как модели Максвелла и Кельвина"--~Фойхта~\cite{kud:bib14}. Эти модели отличаются лишь различными формулами для коэффициентов релаксации. 

Их физический смысл в том, что в них учитывается временная зависимость напряжений и деформаций 

и~их взаимное влияние друг на друга. 

Совпадение моделей \eqref{kud:eq10}, \eqref{kud:eq11} с~моделями

Максвелла, Кельвина"--~Фойхта (с~точностью до постоянных) свидетельствует об использовании одинаковых исходных принципов.



\smallskip



\Section{Математические модели теплопроводности с учетом пространственно"=временной нелокальности}

Рассмотрим последовательность построения локально"=неравновесных моделей теплопроводности на основе формулы \eqref{kud:eq6} для теплового потока в случае бесконечной пластины с толщиной~$\delta$. 

Подставляя \eqref{kud:eq6} в~уравнение теплового баланса \eqref{kud:eq3}, при учете трех членов правой и левой

части соотношения \eqref{kud:eq6} получим уравнение теплопроводности для бесконечной пластины в безразмерном виде

[\citen{kud:bib9}--\citen{kud:bib12}]:

\begin{equation}

\label{kud:eq13}

\frac{\partial  \Theta}{\partial \operatorname{\sf Fo}}+\operatorname{\sf Fo}_1

\frac{\partial ^2 \Theta}{\partial \operatorname{\sf Fo}^2}+\operatorname{\sf Fo}_2

\frac{\partial ^3 \Theta}{\partial \operatorname{\sf Fo}^3}=\frac{\partial

^2 \Theta}{\partial \xi^2}+\operatorname{\sf Fo}_1

\frac{\partial^3 \Theta}{\partial \xi^2\partial

\operatorname{\sf Fo}}+\operatorname{\sf Fo}_2 \frac{\partial ^4 \Theta}{\partial

 \xi^2 \partial \operatorname{\sf Fo}^2},

\end{equation}

где 

$$ 

\Theta=\frac{T-T_{w}}{T_0 -T_{w}},

\quad \xi=\frac{x}{\delta}, \quad 

\operatorname{\sf Fo}=\frac{at}{\delta ^2},

\quad 

\operatorname{\sf Fo}_1 =\frac{a \tau_1}{\delta^2}, 

\quad 

\operatorname{\sf Fo}_2 =\frac{a^2 \tau_2^2}{\delta^4};

$$

%{\color{red}

$T$ "--- температура, $T_{w}$ "--- температура стенки пластины, $T_{0}$ "--- начальная температура, $t$ "--- время,

$a$ "--- коэффициент температуропроводности, 

$\delta $ "--- толщина пластины, 

$\tau_1$ "--- коэффициент релаксации теплового потока,

$\tau_2$ "--- коэффициент релаксации скалярной величины градиента температуры;

 $\Theta$ "--- безразмерная температура, $\operatorname{\sf Fo}$ "--- безразмерное время,

$\xi$ "--- безразмерная координата, 

%$\operatorname{\sf Fo}_1 $ "--- безразмерный коэффициент релаксации теплового потока, 

$\operatorname{\sf Fo}_1 $ "--- безразмерный коэффициент релаксации скалярной величины градиента температуры, 

$\operatorname{\sf Fo}_2 $ "--- безразмерный коэффициент релаксации скалярной величины градиента температуры.

%}







Если положить $\operatorname{\sf Fo}_1 =\operatorname{\sf Fo}_2 =0$ (при $\tau_1=\tau_2=0$), то уравнение \eqref{kud:eq13}

приводится к классическому параболическому уравнению теплопроводности.



Точное аналитическое решение уравнения \eqref{kud:eq13} для бесконечной пластины при

граничных условиях первого рода

$$

 \Theta( 0,\operatorname{\sf Fo})=0, \quad  \Theta( 1,\operatorname{\sf Fo})=0

$$

и начальном условии

$$

 \Theta( \xi,0)=1

$$

имеет вид

\begin{multline}

\label{kud:eq14}

 \Theta( \xi,\operatorname{\sf Fo})=

 \sum\limits_{k=1}^\infty 

\bigl[ D_{1k} \exp (z_{1k} \operatorname{\sf Fo})+D_{2k} \exp (z_{2k} \operatorname{\sf Fo})+\\[-3mm]

+D_{3k} \exp (z_{3k} \operatorname{\sf Fo}) \bigr]

\cos \Bigl(

\frac{2k-1}{2}(1-2 \xi)\pi \Bigr) ,

\end{multline}

где $z_{1k} $, $z_{2k} $, $z_{3k} $ "--- корни характеристического уравнения;

$D_{1k} $, $D_{2k} $, $D_{3k} $ "--- константы интегрирования, 

определяемые из начальных условий краевой задачи.



Результаты расчетов по формуле \eqref{kud:eq14} %{\color{red} (при каких граничных (начальных) условиях получены расчеты?)} 

даны на рис.~\ref{kud:pic2}. 

Их анализ позволяет заключить, что при малых значениях чисел $\operatorname{\sf Fo}_1 $ и $\operatorname{\sf Fo}_2 $

получаемые решения в диапазоне чисел $10^{-6}\le \operatorname{\sf Fo}<\infty $ полностью совпадают с решением аналогичной задачи для параболического уравнения нестационарной теплопроводности.



При $\operatorname{\sf Fo}<10^{-6}$ распределение температуры во времени имеет следующие особенности. 

При значениях величин $\operatorname{\sf Fo}_1 =10^{-6}$, $\operatorname{\sf Fo}_2 =10^{-12}$, например, при

$\operatorname{\sf Fo}=10^{-4}$ температура тела в~бесконечно малых окрестностях точек $ \xi=0$ и $ \xi=1$ скачкообразно изменяется от величины $ \Theta(0, 10^{-4})= \Theta(1, 10^{-4})=0$ непосредственно на границах (выполнение граничных условий первого рода) до $ \Theta(0.0001, 10^{-4})= \Theta(0.9999, 10^{-4})=0.91$ в точках $\xi=0.0001$ и $\xi=0.9999$, 

т.~е. на некотором малом расстоянии от точек $\xi=0$ и $ \xi=1$ (см.~рис.~\ref{kud:pic2} сверху).





\begin{figure}[b!]\centering \footnotesize

\includegraphics[width=\textwidth]{kudinov2}\\



\vspace{-5cm}



\begin{tabular}{p{6.1cm}p{6.7cm}}

&

\caption{\label{kud:pic2} 

Распределение температуры $\Theta$  для различных чисел $\sf Fo$  (значения чисел "--- метки у кривых) в~пластине (см. формулу (\ref{kud:eq14})) для $\operatorname{\sf Fo}_1 =10^{-6}$, $\operatorname{\sf Fo}_2 =10^{- 12}$, $n=100\,000$  (сверху) и  $\operatorname{\sf Fo}_1 =5$, $\operatorname{\sf Fo}_2 =25$, $n=1\,000$ (снизу слева); $n$ "--- число членов ряда (\ref{kud:eq14})}





[Figure~\ref{kud:pic2}. The temperature distribution in the plate (see Eq.~(\ref{kud:eq14})) for different Fourier numbers($\sf Fo$)  (the~values of the Fourier numbers are the labels of the curves) for  $\operatorname{\sf Fo}_1 =10^{-6}$, $\operatorname{\sf Fo}_2 =10^{- 12}$, $n=100\,000$ (top) and  ${\operatorname{\sf Fo}_1 =5}$, $\operatorname{\sf Fo}_2 =25$, $n=1\,000$ (bottom left); $n$ is the  \centerline{number of terms in the series (\ref{kud:eq14})]}



\end{tabular}



\vspace{-.7cm}





\end{figure}







При увеличении чисел $\operatorname{\sf Fo}_1 $, $\operatorname{\sf Fo}_2$ температурные кривые становятся более пологими и~при

$\operatorname{\sf Fo}_1>5 $ "--- практически параллельными оси~$\xi$.

Скачок температуры на стенках (в~окрестностях точек $\xi=0$ и~$\xi=1$) в~этом случае имеет место практически для всех чисел $\operatorname{\sf Fo}$ вплоть до наступления стационарного состояния (см.~рис.~\ref{kud:pic2} снизу слева).



Изменение безразмерного теплового потока $\bar {q}=q\delta /(\lambda (T_0

-T_{w} ))$, найденное по формулам \eqref{kud:eq1}, \eqref{kud:eq2}, \eqref{kud:eq9}, приведено на рис.~\ref{kud:pic3}.



\begin{figure}[t!]

\footnotesize



\begin{tabular}{p{6.1cm}p{6.7cm}}

\includegraphics[width=0.47\textwidth]{kudinov3}&

\vspace{-4.75cm}

\caption{\label{kud:pic3} Изменение теплового потока $\bar q$ на стороне пластины $\xi=1$ при $\operatorname{\sf Fo}_1 =0.3$ по формуле~(\ref{kud:eq1}) (кривая {\sl 1}), по формуле (\ref{kud:eq9}) (кривая~{\sl 2}) и по формуле (\ref{kud:eq2}) (кривая {\sl 3})}

\smallskip



[Figure~\ref{kud:pic3}. The change in the heat flux ($ \bar q $) on the side  $ \xi = 1 $ of the plate when $\operatorname{\sf Fo}_1 =0.3$ by  Eq.~(\ref{kud:eq1}) (curve {\sl 1}), by~Eq.~(\ref{kud:eq9}) (curve~{\sl 2}), and by Eq.~(\ref{kud:eq2}) \centerline{(curve~{\sl 3})]}

\end{tabular}



\vspace{-8mm}





\end{figure}



Анализ результатов расчетов позволяет сделать следующие выводы. 

Тепловой поток, определяемый по формуле \eqref{kud:eq1}, непосредственно на стороне пластины $\xi=1$ при

$\operatorname{\sf Fo}=0$ равен бесконечно большой величине (кривая {\sl 1} на рис.~\ref{kud:pic3}). 

С~увеличением времени величина $\bar {q}(\operatorname{\sf Fo})$ уменьшается по экспоненциальной зависимости и~при $\operatorname{\sf Fo}\approx 3$ становится практически равной нулю.



Тепловой поток, определяемый по формуле, представляющей решение уравнения \eqref{kud:eq9}, при $\xi =1$ и $\operatorname{\sf Fo}=0$ равен нулю. 

С~увеличением времени он возрастает и при $\operatorname{\sf Fo}=\operatorname{\sf Fo}^{\ast }\approx 0.072$ достигает максимального значения $\bar {q}\approx 0.37$ (кривая {\sl 2} на рис.~\ref{kud:pic3}).



Изменение теплового потока, определяемого формулой \eqref{kud:eq2}, иллюстрируется кривой {\sl 3}

на рис.~\ref{kud:pic3}. 

В~данном случае тепловой поток на стороне $\xi =1$, имея нулевое значение при $\operatorname{\sf Fo}=0$, с~увеличением числа $\operatorname{\sf Fo}$ возрастает до некоторого максимального значения, равного $\bar {q} \approx 1.83$ (при $\operatorname{\sf Fo}= \operatorname{\sf Fo}'\approx 0.08)$. 

При дальнейшем увеличении числа $\operatorname{\sf Fo}$ тепловой поток изменяется скачкообразно, принимая в~некоторых диапазонах числа $\operatorname{\sf Fo}$ отрицательные значения. 

Такое изменение теплового потока приводит к~появлению скачков температуры и~к~их

отрицательным значениям.



Обнаруженный из анализа точного аналитического решения \eqref{kud:eq14} факт задержки во

времени принятия граничного условия первого рода свидетельствует о том, что

ввиду сопротивления тела, оказываемого процессу проникновения теплоты, его

мгновенный прогрев на границе невозможен ни при каких условиях теплообмена с

окружающей средой. Следовательно, коэффициент теплоотдачи на стенке зависит

не только от условий теплообмена (скорость среды, вязкость и проч.), но и от

физических свойств тела, и он, во"=первых, является переменной во времени

величиной и, во"=вторых, он не может превысить некоторой предельной для

каждого конкретного случая величины. Конкретный пример определения

коэффициента теплоотдачи для стальной пластины ($ \lambda =49$

Вт/(м${}\cdot{}$К)) толщиной $ \delta =0.01$~м при $\operatorname{\sf Fo}_ 1 =0.3$,

$\operatorname{\sf Fo}_2 =0$, который оказался равным $\alpha=4005$~Вт/(м$^{2}\cdot{}$К), приведен в~\cite{kud:bib12}. Таким образом, механизм теплопроводности с~учетом релаксационных свойств материалов ограничивает количество тепла,

отводимого от системы. 

Этот факт необходимо учитывать при рассмотрении всех быстропротекающих процессов "--- при решении задач теплового воспламенения (теплового взрыва), задач динамической термоупругости, где в~условиях теплового удара классическая теория приводит к~бесконечным динамическим напряжениям, при лазерной обработке и~упрочнении материалов и~во многих

других процессах.



\smallskip







\Section{Математическая модель колебаний давления в несжимаемой жидкости

с учетом пространственно"=временной нелокальности} \linebreak

С~учетом соотношения \eqref{kud:eq9} авторами работы~\cite{kud:bib10} было получено дифференциальное

уравнение и сформулирована краевая задача, описывающая колебание давления упругой жидкости в~трубопроводе в~условиях гидравлического удара. 

Математическая постановка задачи в данном случае имеет вид

\begin{multline}

\label{kud:eq15}

\operatorname{\sf Fo}_r \frac{\partial  \Theta( \xi,\operatorname{\sf Fo})}{\partial \operatorname{\sf Fo}}+\frac{\partial^2 \Theta( \xi,\operatorname{\sf Fo})}{\partial

\operatorname{\sf Fo}^{2}}+\operatorname{\sf Fo}_1 \frac{\partial^3 \Theta( \xi,\operatorname{\sf Fo})}{\partial \operatorname{\sf Fo}^3}=\\

=\frac{\partial^2 \Theta( \xi,\operatorname{\sf Fo})}{\partial

 \xi^{2}}+\operatorname{\sf Fo}_2 \frac{\partial^3 \Theta( \xi,\operatorname{\sf Fo})}{\partial  \xi^2\partial

\operatorname{\sf Fo} }, \quad \operatorname{\sf Fo}\geq 0,\quad 0 \leq \xi \leq 1;

\end{multline}

\begin{gather}

\label{kud:eq16}

\Theta(\xi , {\sf Fo})\big|_{\mathop{\sf Fo}=0}=1,

\quad

\frac{\partial  \Theta( \xi, \mathop{\sf Fo} )}{\partial \operatorname{\sf Fo}}\Big|_{\mathop{\sf Fo}=0}=0,

\quad

\frac{\partial^2 \Theta( \xi, \mathop{\sf Fo} )}{\partial \operatorname{\sf Fo}^2}\Big|_{\mathop{\sf Fo}= 0}=0;\\

\label{kud:eq17}

\frac{\partial  \Theta(\xi,\operatorname{\sf Fo})}{\partial  \xi}\Big|_{\xi=0}=0,

\quad

\Theta(\xi,\operatorname{\sf Fo})\big|_{\xi=1}=0.

\end{gather}

В данной модельной задаче используются следующие обозначения:

$$ \Theta=\frac{p-p_1 }{p_0 -p_1 },\quad 

\operatorname{\sf Fo}=\frac{c_st}l, \quad 

\xi=1-\frac{x}l, $$

$$

\operatorname{\sf Fo}_ 1 =\frac{c_s \tau_1 }l, \quad

\operatorname{\sf Fo}_{2} =\frac{c_s \tau_2 }{l},\quad 

\operatorname{\sf Fo}_r =\frac{2 b l}{c_s},\quad 

2b=\frac{32\nu}{d^2};

$$

$p$ "--- давление, $p_0 $ "--- начальное давление, $p_1 $ "--- давление на входе в~трубопровод, 

$x$ "--- координата, $l$ --- длина трубы, $c_s$ "--- скорость звука в жидкости; 

$\nu $ "--- вязкость жидкости, $d$ "--- диаметр трубы; 

$\tau_1$ "--- коэффициент релаксации касательного напряжения, 

$\tau_2$ "--- коэффициент релаксации градиента скорости;

$\Theta$ "--- безразмерное давление, 

$\operatorname{\sf Fo}$ "--- безразмерное время,

$\xi$ --- безразмерная координата,

$\operatorname{\sf Fo}_ 1 $ "--- безразмерный коэффициент релаксации касательного напряжения, 

$\operatorname{\sf Fo}_ {2}$ "--- безразмерный коэффициент релаксации градиента скорости,

$\operatorname{\sf Fo}_r$ "--- безразмерный коэффициент сопротивления среды.





Краевая задача \eqref{kud:eq15}--\eqref{kud:eq17} описывает распространение скачка давления в~трубопроводе с~неподвижной в~исходном состоянии жидкостью для случая, когда в сечении $\xi =1$ произошло скачкообразное изменение давления, а~сечение $\xi =0$ перекрыто (скорость течения жидкости равна нулю). 



Точное аналитическое решение задачи \eqref{kud:eq15}--\eqref{kud:eq17} находится в виде

\begin{multline}

\label{kud:eq18}

 \Theta( \xi,\operatorname{\sf Fo})=

 \sum\limits_{k=1}^\infty 

 \bigl[C_{1k} \exp \left( z_{1k}

\operatorname{\sf Fo} \right)+C_{2k} \exp \left( z_{2k} \operatorname{\sf Fo} \right) + \\ +C_{3k} \exp

\left( z_{3k} \operatorname{\sf Fo} \right) \bigr] 

\cos \Bigl( \frac{2k-1}2 \pi

 \xi \Bigr),

\end{multline}

где $C_{1k} $, $C_{2k} $, $C_{3k} $ --- константы интегрирования,

определяемые из начальных условий \eqref{kud:eq16}; $z_{1k} $, $z_{2k} $, $z_{3k} $ ---

корни характеристического уравнения.



Результаты расчетов по формуле \eqref{kud:eq18} приведены на рис.~\ref{kud:pic4} (для $ \xi=0$). 

Здесь даны результаты решения задачи \eqref{kud:eq15}--\eqref{kud:eq17} в квазистационарной

постановке (при $\operatorname{\sf Fo}_1 =\operatorname{\sf Fo}_2 =0$), а также для случая, когда

учитывается конечная скорость распространения потенциалов исследуемой среды

($\operatorname{\sf Fo}_ 1 \ne 0$, $\operatorname{\sf Fo}_2 \ne 0$), то есть с учетом

влияния нестационарности процесса. На этих рисунках представлены также

результаты расчетов, выполненных автором работы~\cite{kud:bib15}, и результаты

экспериментальных исследований гидравлического удара для сильновязкой

жидкости (масло) при ламинарном режиме течения~\cite{kud:bib16}. 

Теоретические исследования в~\cite{kud:bib15} выполнялись по двум методикам "--- с~учетом влияния

нестационарности процесса на силу трения и без ее учета (в квазистатической

постановке).





\begin{figure}[t!]\centering \footnotesize

\includegraphics[scale=0.7]{kudinov4}\\

\caption{Распределение давления в трубопроводе при гидравлическом ударе ($ \xi=0$), когда $\operatorname{\sf Fo}_r =0.0538$ и $\operatorname{\sf Fo}_1 =\operatorname{\sf Fo}_2 =0$ (слева) и $\operatorname{\sf Fo}_r =0.0538$,

$\operatorname{\sf Fo}_1 = 0.4$, $\operatorname{\sf Fo}_2 =0.043$ (справа): 

{\sl 1}~--- расчет по формуле (\ref{kud:eq18}) при $n=100$; 

{\sl 2}~--- по методу~\cite{kud:bib15};

{\sl 3}~--- эксперимент~\cite{kud:bib16}

\label{kud:pic4}}



\smallskip



[Figure~\ref{kud:pic4}. The pressure distribution in a pipeline under the conditions of a hydraulic shock

($ \xi=0$) when $\operatorname{\sf Fo}_r =0.0538$, $\operatorname{\sf Fo}_1 =\operatorname{\sf Fo}_2 =0$ (left), and $\operatorname{\sf Fo}_r =0.0538$,

$\operatorname{\sf Fo}_1 = 0.4$, $\operatorname{\sf Fo}_2 =0.043$ (right): 

curve {\sl 1}~--- calculation by Eq.~(\ref{kud:eq18}) when $n=100$; 

curve {\sl 2}~--- by method~\cite{kud:bib15};

curve {\sl 3}~--- experimental data~\cite{kud:bib16}]



\vspace{-6mm}



\end{figure}











Анализ результатов расчетов по формуле \eqref{kud:eq18} позволяет сделать вывод о~том, что

решение задачи \eqref{kud:eq15}--\eqref{kud:eq17} в~квазистатической постановке (при $\operatorname{\sf Fo}_1

=\operatorname{\sf Fo}_2 =0$) незначительно отличается от решения аналогичной задачи,

полученного в~\cite{kud:bib15} с~использованием метода Бернулли"--~Фурье. 

Однако результаты двух теоретических методов значительно отличаются от

экспериментальных данных. 

Учет нестационарности градиента скорости и~касательного напряжения приводит к~существенному приближению расчетных

данных к~результатам эксперимента. 

На основе экспериментальных данных, используя соотношение \eqref{kud:eq18}, путем решения обратной задачи были найдены коэффициенты релаксации: $\tau _1 =0.0108$~с, $\tau _2 =0.0116$~с.





Отметим, что в классическом гиперболическом уравнении, описывающем

распространение скачка давления при гидравлическом ударе, когда $\operatorname{\sf Fo}_1

\hm=\operatorname{\sf Fo}_2 =0$, зависимости между значениями скорости, плотности и~давления потока жидкости принимаются осредненными по его сечению. 

Кроме этого, они содержат среднюю во времени величину касательного напряжения на

стенке трубопровода, то есть используется гипотеза квазистационарности, 

в~соответствии с~которой считается, что закон Ньютона для касательного

напряжения выполняется и~при нестационарном процессе изменения скорости.

Предполагается также, что соответствующие соотношения сохраняются и~для

коэффициента гидравлического сопротивления. 

Строгое обоснование этого допущения отсутствует, поэтому его справедливость может быть подтверждена

лишь путем сравнения результатов расчета с~экспериментальными исследованиями.



Полученные двумя различными методами результаты позволяют заключить, что

гипотеза квазистационарности, вообще говоря, неверна, что подтверждается

приведенными выше сравнениями решений квазистационарной и нестационарной

задач с~результатами натурного эксперимента. 

Из двух рассмотренных выше направлений учета нестационарности касательных напряжений первое связано с~рассмотрением дифференциальных уравнений Навье"--~Стокса, описывающих изменение местных (локальных "--- неосредненных) величин~\cite{kud:bib15}. 

Второе направление, рассмотренное в данной работе, связано с~введением временн\'{o}й зависимости между градиентом скорости и~касательным напряжением в формуле \eqref{kud:eq10} закона Ньютона. 

В данном случае устраняется проблема бесконечной скорости распространения потенциалов исследуемого поля,

а~также проблема бесконечных величин касательных напряжений на стенке трубопровода при

$\operatorname{\sf Fo}\to 0$.



Анализ безразмерного касательного напряжения $ {\bar \tau}$, определяемого на основе формулы \eqref{kud:eq10}, которая при $ \tau_1 = \tau_2 =0$ приводится к формуле закона Ньютона, позволяет заключить, что $\bar \tau \to \infty$ при $\operatorname{\sf Fo}\to 0$ (см. рис.~\ref{kud:pic5}). 

При $ \tau_1 $ и~$ \tau_2 $, отличных от нуля, касательное напряжение $ {\bar \tau}$, имея в начальный момент времени нулевое значение, с~увеличением времени возрастает, принимая при некотором $\operatorname{\sf Fo}=\operatorname{\sf Fo}^\ast $ максимальное значение $ {\bar \tau}= {\bar \tau^\ast}$. 

При дальнейшем возрастании времени $ {\bar \tau}$ уменьшается и~при $\operatorname{\sf Fo} \to \infty$ $\bar \tau \to 0$. 

Полученные результаты применительно к решению уравнения \eqref{kud:eq10} позволяют

сделать следующие выводы: 

на стенке трубопровода касательное напряжение  $\bar \tau \to 0$ при $\operatorname{\sf Fo} \to 0$  и~не может превысить некоторой максимальной для данных конкретных условий (определяемых физическими свойствами среды)

величины; применение формулы \eqref{kud:eq10} позволяет избежать бесконечных значений касательных напряжений.



\begin{figure}\centering \footnotesize







\begin{tabular}{p{6.2cm}p{6.4cm}}

\includegraphics[scale=0.8]{kudinov5}~&

\vspace{-2.75cm}

\caption{Изменение касательного напряжения на стенке трубопровода при ${\operatorname{\sf Fo}_r =2al/c=0.3}$; 

решение уравнения (\ref{kud:eq10}) при ${ \tau_1 = \tau_2 =0}$ (линия~{\sl 1}) и при 

$ \tau_1 =0.0108$, $\tau_2 =0.0116$ (линия {\sl 2})

\label{kud:pic5}}

\end{tabular}





[Figure~\ref{kud:pic5}. 

The change of the tangential stress ($\bar \tau$)  distribution on the pipeline wall, which was found as solution of Eq.~(\ref{kud:eq10})

when  $ \tau_1 = \tau_2 =0$ (curve~{\sl 1}), and~$ \tau_1 =0.0108$, $\tau_2 =0.0116$ (curve~{\sl 2}); here $\operatorname{\sf Fo}_r =2al/c=0.3$]



\vspace{-6mm}



\end{figure}







Изменение давления в случае $\operatorname{\sf Fo}_1 =\operatorname{\sf Fo}_2 =0$ (квазистационарный

режим) приведено на левом рис.~\ref{kud:pic4} (линия {\sl 1}). 

Анализ результатов позволяет сделать заключение о скачкообразном изменении давления в любой точке трубы

во времени, свидетельствующем о бесконечной скорости передачи импульса.

Отметим, что скорость перемещения волны при этом вполне конечна и равна скорости звука в жидкости. 

Бесконечной является скорость передачи импульса, заложенная в классической (нерелаксированной) формуле закона Ньютона для касательного напряжения. 

Полученные результаты позволяют сделать вывод об ограниченности математических моделей колебаний упругих жидкостей, в~которых не учитываются их релаксационные свойства, ввиду существенного отклонения получаемых данных от результатов экспериментальных исследований (см. рис.~\ref{kud:pic4} слева). 

Учет релаксационных свойств позволяет через модификацию закона Ньютона для касательного напряжения учесть нестационарность изменения скорости и~касательного напряжения.



\smallskip







\Section{Математическая модель теплообмена в движущейся жидкости

с учетом ее релаксационных свойств}

В предложенной выше математической модели \eqref{kud:eq15}--\eqref{kud:eq17} исследуется изменение давления в~жидкости в условиях гидравлического удара. 

Рассмотрим вывод уравнения применительно к~нестационарному теплообмену при ламинарном

течении вязкой несжимаемой жидкости в~плоскопараллельном канале  толщиной $2\delta$ (см. рис.~\ref{kud:pic6}). 

Формулы для теплового потока в~направлении координатных осей $x$ и~$y$ по аналогии с \eqref{kud:eq9} будут иметь вид~\cite{kud:bib17}

\begin{equation}

\label{kud:eq19}

\begin{aligned}

 q_x =&- \lambda\cfrac{\partial T}{\partial x}- \tau_1

\frac{\partial q_x }{\partial t}- \lambda \tau _2 \frac{\partial

^2T}{\partial x\partial t}+ \rho \omega _x h; \\

 q_y =&- \lambda \frac{\partial T}{\partial y}- \tau_1

\frac{\partial q_y }{\partial t}- \lambda \tau_2 \frac{\partial

^2T}{\partial y\partial t}+ \rho \omega_y h,

 \end{aligned}

\end{equation}

где 

$q_x $, $q_y $ "--- тепловые потоки по осям $x$ и $y$; 

$\omega_x $, $\omega_y$ "--- скорости течения жидкости в направлении координатных осей $x$ и $y$; 

$h$ "--- теплосодержание (энтальпия жидкости);

$\lambda$ "--- коэффициент теплопроводности;

$x$, $y$ "--- продольная и поперечная координаты плоскопараллельного канала;

$ \tau_ 1$, $\tau_2$ "--- коэффициенты релаксации соответственно теплового потока и градиента температуры.



\begin{figure}[t!]



\footnotesize \centering



\begin{tabular}{p{5.2cm}p{7.4cm}}

\includegraphics[scale=0.8]{kudinov6}~&

\vspace{-3.75cm}

\caption{Схема теплообмена при ламинарном течении жидкости в плоскопараллельном канале

\label{kud:pic6}}

\\

& \vspace{-2.85cm}

[Figure~\ref{kud:pic6}. 

A heat exchange scheme for laminar flow of liquid in the plane-parallel channel; 

$T_{i}$~is the initial liquid temperature; $T_0$ is the liquid temperature at the channel inlet;

$T_{w}$ is the liquid temperature on the channel wall;

$2\delta$ is the channel width; 

$\omega_x $ is the fluid \centerline{flow velocity in~the direction of the $x$-axis]}



\end{tabular}











\vspace{-5mm}



\end{figure}







Отличие формул \eqref{kud:eq19} от соотношения \eqref{kud:eq9} состоит в том, что при движении

жидкости теплота переносится не только теплопроводностью, но и конвекцией,

учитываемой последними слагаемыми этих формул.



Уравнение теплового баланса в данном случае имеет вид

\begin{equation}

\label{kud:eq20}

\rho\frac{\partial h}{\partial t}=

c\rho  \frac{\partial T}{\partial t}=

-\Bigl( \frac{\partial q_x

}{\partial x}+\frac{\partial q_y}{\partial y} \Bigr).

\end{equation}





Подставляя \eqref{kud:eq19} в \eqref{kud:eq20}, пренебрегая конвективным переносом теплоты в~поперечном направлении ($\rho \omega_y \frac{\partial h}{\partial y}=0$) и

теплопроводностью "--- в продольном ($ \lambda \frac{\partial^2T}{\partial

x^2}=0$), получаем

\begin{equation}

\label{kud:eq21}

\frac{\partial T}{\partial t}+ {\omega}_x \frac{\partial T}{\partial

x}=a\frac{\partial^2T}{\partial y^2}- \tau_1 \frac{\partial

^2T}{\partial t^2}+a \tau_2 \frac{\partial }{\partial t}\left(

\frac{\partial ^2T}{\partial x^2}+\frac{\partial ^2T}{\partial y^2}

\right),

\end{equation}

где 

$$a= \frac{\lambda}{c\rho},\quad \omega_x=\frac{3}{2} \omega_{av} \Bigl( 1-\frac{y^2}{\delta ^2}\Bigr);$$ $\omega_{av}$ "--- средняя скорость

течения жидкости.







В данной модели будем использовать следующие обозначения:

 $T_{i}$ "--- начальная температура жидкости; $T_0$ "--- температура

жидкости на входе в канал; $T_{w}$ "--- температура жидкости на

стенке; $2\delta$ "--- ширина плоского канала, а~также  следующие  безразмерные переменные и параметры: 

\begin{equation}

\label{kud:eq22}

\begin{aligned}

\Theta=\frac{T-T_{w}}{T_0 -T_{w} };

\quad

\eta =\frac{y}{ \delta };

\quad

\operatorname{\sf Fo}=\frac{at}{\delta^2}; \quad

\operatorname{\sf Fo}_1^* =\frac{a \tau_1}{ \delta^2}; \\

\operatorname{\sf Fo}_2^* =\frac{a \tau_2 }{ \delta^2}; \quad

\xi=\frac{2}{3}\frac{ax}{ \delta^ 2 \omega_{av} }; \quad

\operatorname{\sf Pe}=\Bigl(\frac{2}{3}\frac{a}{\delta \omega_{av} }

\Bigr)^2.

\end{aligned}

\end{equation}

С учетом обозначений \eqref{kud:eq22} уравнение \eqref{kud:eq21} приводится к виду

\begin{equation}

\label{kud:eq23}

\begin{array}{c}

\displaystyle 

\frac{\partial  \Theta}{\partial \operatorname{\sf Fo}}+

\operatorname{\sf Fo}_1 ^* \frac{\partial^2 \Theta}{\partial \operatorname{\sf Fo}^ {2}}+(1- \eta

)\frac{\partial  \Theta}{\partial  \xi} =\frac{\partial

^2 \Theta}{\partial \eta^2}+\operatorname{\sf Fo}_2^* \frac{\partial

}{\partial \operatorname{\sf Fo}}

\Bigl(\frac{\partial^2 \Theta}{\partial  \xi^2} \operatorname{\sf Pe}+\frac{\partial^2 \Theta}{\partial  \eta

^2} \Bigr),

\\[3mm]

\operatorname{\sf Fo}>0, \quad  \xi>0, \quad  0<\eta<1.

\end{array}

\end{equation}



Начальные и граничные условия имеют вид

\begin{equation}

\label{kud:eq24}

\begin{aligned}

&\Theta (\xi ,\eta,0)= \Theta_ {i},

\quad

\frac{\partial  \Theta(\xi,\eta,\operatorname{\sf Fo})}{\partial \operatorname{\sf Fo}}\Big|_{\operatorname{\sf Fo}=0}=0; \\

& \Theta(0, \eta,\operatorname{\sf Fo})=1, \quad 

\frac{\partial  \Theta(\xi, \eta, \operatorname{\sf Fo})}{\partial  \xi}\Big|_{\xi = \infty}=0;

\\

&

\frac{\partial  \Theta( \xi, \eta,\operatorname{\sf Fo})}{\partial \eta}\Big|_{\eta=0}=0, \quad 

\Theta( \xi, 1,\operatorname{\sf Fo})=0,

\end{aligned}

\end{equation}

где $\Theta_{i} =({T_ {i} -T_{w}})/({T_0

-T_{w}})$ "--- безразмерная начальная температура.



Для решения задачи \eqref{kud:eq23}, \eqref{kud:eq24} использовался метод конечных разностей.

Результаты расчетов (для $\Theta_{i} =0.5$) представлены на рис.~\ref{kud:pic7}, \ref{kud:pic8}. 

На рис.~\ref{kud:pic7} даны результаты, полученные при равенстве нулю последних слагаемых правой и~левой частей уравнения \eqref{kud:eq23}, которое в~этом случае приводится к уравнению \eqref{kud:eq4}. 

Анализ результатов расчетов позволяет сделать заключение о~практическом совпадении расчетов численным методом с~точным

аналитическим решением уравнения \eqref{kud:eq4}.



\begin{figure}[b!]\centering \small



\footnotesize \centering



\begin{tabular}{p{5.6cm}p{7cm}}

\includegraphics[scale=0.8]{kudinov7}~~~&

\vspace{-5cm}

\caption{Распределение температуры в плоском канале с~неподвижной жидкостью ($\operatorname{\sf Fo}_1^* =\operatorname{\sf Fo}_ 2^* \hm =10^{-7}$) для различных чисел {\sf Fo} (значения чисел~— метки у кривых): сплошные линии "--- точное решение [\citen{kud:bib9}--\citen{kud:bib13}];  $\tiny\boldsymbol{\circ}$ "--- метод конечных разностей

\label{kud:pic7}}

\\

&\vspace{-3cm}

[Figure~\ref{kud:pic7}.  The temperature distribution in a flat channel  channel with a fixed liquid volume (when $\operatorname{\sf Fo}_1^* =\operatorname{\sf Fo}_ 2^* \hm =10^{-7}$) for different Fourier numbers ({\sf Fo}) (the values of the Fourier numbers are the labels of the curves);

the solid lines are the exact solution [\citen{kud:bib9}--\citen{kud:bib13}] and  lines of points ($\tiny\boldsymbol{\circ}$) are the \centerline{simulation results  by the finite difference method]}



\\



\includegraphics[scale=0.8]{kudinov8}~~~&

\vspace{-5cm}

\caption{Распределение температуры по ширине канала при $ \xi=5\cdot 10^{-5}$, $\operatorname{\sf Fo}_1^* =\operatorname{\sf Fo}_2^* =10^{-7}$, $ {\sf Pe}=10^{-7}$ в диапазоне безразмерного времени $2\cdot 10^{-8}\le \operatorname{\sf Fo}\le 10^{-4}$: сплошные линии "---

точное решение [\citen{kud:bib9}--\citen{kud:bib13}];  $\tiny\boldsymbol{\circ}$ "--- метод конечных разностей}

\label{kud:pic8}

\\

&\vspace{-3cm}

[Figure~\ref{kud:pic8}. The temperature distribution along the channel width (when $ \xi=5\cdot 10^{-5}$, $\operatorname{\sf Fo}_1^* =\operatorname{\sf Fo}_2^* =10^{-7}$, $ {\sf Pe}=10^{-7}$) for  $2\cdot 10^{-8}\le \operatorname{\sf Fo}\le 10^{-4}$ (the values of the Fourier numbers are the labels of the curves);

the solid lines are the exact solution [\citen{kud:bib9}--\citen{kud:bib13}] and  lines of points ($\tiny\boldsymbol{\circ}$) are the simulation \centerline{results  by the finite difference method]}





\end{tabular}



\end{figure}



На рис.~\ref{kud:pic8} дано распределение температуры жидкости при малых значениях

величин поперечной координаты $0.99 \le  \eta \le 1.0$ и чисел Фурье из диапазона

${2\cdot 10^{-8}}\le \operatorname{\sf Fo}\le 10^4$. 

Анализ полученных результатов позволяет отметить интересную особенность теплообмена при задании граничных

условий первого рода (тепловой удар) на стенке в случае, когда учитываются релаксационные свойства среды. 

Так, несмотря на выполнение граничного условия первого рода $ \Theta (\xi,1,\operatorname{\sf Fo})=0$ в~точке 

$ \xi=1$, в бесконечно малой окрестности этой точки, например, $ \xi=0.999$, температура стенки не принимает значение $ \Theta ( \xi,1,\operatorname{\sf Fo})=0$ мгновенно, то есть непосредственно при $\operatorname{\sf Fo}=0$. 

И, в частности, для того чтобы температура стенки приняла значение $ \Theta ( \xi,1,\operatorname{\sf Fo})=0$, требуется некоторый промежуток времени, равный $\operatorname{\sf Fo}=10^{-4}$ (при $\operatorname{\sf Fo}_1^* =\operatorname{\sf Fo}_2^* =10^{-7}$). 

Полученные результаты полностью согласуются с~результатами точного аналитического решения вида \eqref{kud:eq14} [\citen{kud:bib9}--\citen{kud:bib13}].



Отметим, что при некоторых больших значениях числа Фурье теплообмен стабилизируется и наступает квазистационарный режим, когда температура жидкости в любой точке канала принимает неизменное во времени значение.





\smallskip



\Section{Локально"=неравновесные модели процессов колебаний упругих твердых тел}

Используя модифицированную формулу закона Гука вида \eqref{kud:eq11}, рассмотрим последовательность построения локально"=неравновесной модели колебаний упругого стержня длины $\delta$ с~учетом сопротивления, оказываемого материалом процессу

колебаний. 

Уравнение равновесия (движения) в данном случае имеет вид

\begin{equation}

\label{kud:eq25}

\rho \frac{\partial ^2U}{\partial t^2}=\frac{\partial  \sigma}{\partial x}-\gamma\rho \frac{\partial U}{\partial t},

\end{equation}

где 

$ {\gamma }$ "--- коэффициент сопротивления;

$\rho $ "--- плотность.



Отметим, что второе слагаемое правой части уравнения \eqref{kud:eq25} представляет силу

сопротивления среды, приходящуюся на единицу объема.



Подставляя \eqref{kud:eq11} в \eqref{kud:eq25}, с учетом того, что 

$$ \frac{\partial  \sigma}{\partial x}= \rho \frac{\partial ^2U}{\partial t^2},$$ находим

\begin{equation}

\label{kud:eq26}

\tau_1^c \frac{\partial ^3U}{\partial t^3}+\frac{\partial

^2U}{\partial t^2}=e^2_s\frac{\partial ^2U}{\partial x^2}+e^2_s \tau_2^c

\frac{\partial ^3U}{\partial x^2\partial t}- \gamma \frac{\partial U}{\partial t},

\end{equation}

где 

$e_s=\sqrt {E/\rho }$ "--- скорость распространения волны (скорость звука в~материале стержня).



Очевидно, что при $ \tau_1^c = \tau_2^c = \gamma =0$

уравнение \eqref{kud:eq26} приводится к волновому уравнению незатухающих колебаний:

\begin{equation}

\label{kud:eq27}

\frac{\partial^2U}{\partial t^2}=e_s\frac{\partial ^2U}{\partial x^2}.

\end{equation}



Найдем точное аналитическое решение краевой задачи о колебаниях стержня,

один торец которого жестко закреплен, а второй   находится в свободном

состоянии. В начальный момент времени стержень деформирован по линейному

закону, согласно которому максимальное перемещение имеет свободный торец

стержня. 



Воспользовавшись безразмерными переменными

\begin{equation}

\label{kud:eq28}

\Theta=\frac{U}{U_0}, \quad  

\xi=\frac{x}{ \delta},\quad 

\operatorname{\sf Fo}=\frac{e_st}{ \delta},\quad 

\operatorname{\sf Fo}_1^c =\frac{e_s \tau_ 1^c }{\delta},\quad 

\operatorname{\sf Fo}_2^c =\frac{e_s \tau_2^c }{\delta},\quad 

\operatorname{\sf Fo}_3^c =\frac{\delta \gamma}{e_s},

\end{equation}

рассматриваемую задачу запишем в виде

\begin{equation}

\label{kud:eq29}

\begin{array}{c}

\displaystyle

\operatorname{\sf Fo}_3^c \frac{\partial  \Theta( \xi,\operatorname{\sf Fo})}{\partial

\operatorname{\sf Fo}}+

\operatorname{\sf Fo}_1^c \frac{\partial ^3 \Theta( \xi,\operatorname{\sf Fo})}{\partial \operatorname{\sf Fo}^3}+

\frac{\partial ^2 \Theta( \xi,\operatorname{\sf Fo})}{\partial \operatorname{\sf Fo}^2}

=\frac{\partial^2 \Theta( \xi,\operatorname{\sf Fo})}{\partial  \xi^2}+

\operatorname{\sf Fo}_2^c \frac{\partial^3 \Theta( \xi,\operatorname{\sf Fo})}{\partial  \xi^2\partial \operatorname{\sf Fo}}, \\[2mm]

\operatorname{\sf Fo}>0, \quad  0< \xi<1;

\end{array}

\end{equation}

\begin{gather}

\label{kud:eq30}

\Theta( \xi,0)=1,

\quad

\frac{\partial  \Theta( \xi,\operatorname{\sf Fo})}{\partial \operatorname{\sf Fo}}\Big|_{\operatorname{\sf Fo}=0}=0,

\quad

\frac{\partial^2 \Theta( \xi, \operatorname{\sf Fo})}{\partial \operatorname{\sf Fo}^2}\Big|_{\operatorname{\sf Fo}=0}=0;\\

\label{kud:eq31}

\frac{\partial  \Theta(\xi,\operatorname{\sf Fo})}{\partial  \xi}\Big|_{\xi=0}=0,

\quad

\Theta(1,\operatorname{\sf Fo})=0.

\end{gather}

Здесь $ \Theta$ "--- безразмерное перемещение; $ \xi $ "--- безразмерная координата; $\operatorname{\sf Fo}$ "--- число Фурье (безразмерное время); $U_0 $ "--- начальное перемещение стержня; $\operatorname{\sf Fo}_1^c$, $\operatorname{\sf Fo}_2^c$ "--- безразмерные

коэффициенты релаксации; $\operatorname{\sf Fo}_3^c$ "--- безразмерный коэффициент

сопротивления среды; $ \delta$ "--- длина стержня.



Точное аналитическое решение задачи \eqref{kud:eq29}--\eqref{kud:eq31}, найденное методом

разделения переменных, записывается в виде

\begin{multline}

\label{kud:eq32}

\Theta( \xi,\operatorname{\sf Fo})=\sum\limits_{k=1}^\infty \Bigl\{

\exp (\gamma_k \operatorname{\sf Fo})

\bigl[ B_{1k} \cos ( \beta_k \operatorname{\sf Fo})-B_{2k}  \sin ( \beta _k \operatorname{\sf Fo}) \bigr]+\\

+  B_{3k} \exp (z_{3k} \operatorname{\sf Fo}) \Bigr\} \cos \Bigl(\frac{2k-1}2 \pi \xi \Bigr),

\end{multline}

где $ {\gamma }_k $, $\beta_k$, $z_{3k} $ "--- корни

характеристического уравнения; $B_{1k}$, $B_{2k}$, $B_{3k}$ "--- константы

интегрирования, определяемые из начальных условий \eqref{kud:eq30}.



На рис.~\ref{kud:pic9}, \ref{kud:pic10} приведены результаты расчетов перемещений по формуле \eqref{kud:eq32}

при $\operatorname{\sf Fo}_ 1^c =\operatorname{\sf Fo}_2^c =0.1$, $\operatorname{\sf Fo}_{3}^c =0.5$

(рис.~\ref{kud:pic9}) и 

$\operatorname{\sf Fo}_ 1^c =\operatorname{\sf Fo}_2^c =0.1$, 

$\operatorname{\sf Fo}_3^c=30$ (рис.~\ref{kud:pic10}). 

Из их анализа можно заключить, что с~увеличением числа $\operatorname{\sf Fo}_3 ^c$ (при неизменных $\operatorname{\sf Fo}_1^c$ и $\operatorname{\sf Fo}_2^c$) колебания смещаются в область положительных значений перемещений и~при некоторых больших значениях величины $\operatorname{\sf Fo}_3^c$

($ \operatorname{\sf Fo}_3^c>5$) возврат стержня в исходное состояние происходит практически без колебаний (критическое затухание (рис.~\ref{kud:pic10})).





\begin{figure}[t!]



\footnotesize \centering



\begin{tabular}{p{7.7cm}p{4.8cm}}

\includegraphics[scale=0.8]{kudinov9}~~~&

\vspace{-5.25cm}

\caption{Изменение перемещений стержня во времени для различных $\xi$ при $\operatorname{\sf Fo}_ 1^c \hm =\operatorname{\sf Fo}_2^c =0.1$, $\operatorname{\sf Fo}_ 3^c =0.5$, $n=100$; $n$ "--- число членов ряда (\ref{kud:eq32})

\label{kud:pic9}}

\\

&\vspace{-3.4cm}

[Figure~\ref{kud:pic9}. The change of the rod displacement over time for various~$\xi$ when $\operatorname{\sf Fo}_ 1^c \hm =\operatorname{\sf Fo}_2^c =0.1$, $\operatorname{\sf Fo}_ 3^c =0.5$, $n=100$; 

$n$ is the number of terms in the series (\ref{kud:eq32})]

\\



~~\,\,\includegraphics[scale=0.8]{kudinov10}~~~&

\vspace{-5.25cm}

\caption{Изменение перемещений стержня во времени для различных $\xi$ при $\operatorname{\sf Fo}_ 1^c \hm =\operatorname{\sf Fo}_2^c =0.1$, $\operatorname{\sf Fo}_ 3^c =30$, $n=100$; $n$ "--- число членов ряда (\ref{kud:eq32})

\label{kud:pic10}}

\\

&\vspace{-3.4cm}

[Figure~\ref{kud:pic10}. The change of the rod displacement over time for various~$\xi$ when $\operatorname{\sf Fo}_ 1^c \hm =\operatorname{\sf Fo}_2^c =0.1$, $\operatorname{\sf Fo}_ 3^c =30$, $n=100$; 

$n$ is the number of terms in the series (\ref{kud:eq32})]

\\



\includegraphics[scale=0.8]{kudinov11}~~~&

\vspace{-5.25cm}

\caption{Изменение перемещений стержня во времени для различных $\xi$ при $\operatorname{\sf Fo}_ 1^c \hm ={\operatorname{\sf Fo}_2^c =0.1}$, $\operatorname{\sf Fo}_ 3^c =0.5$, $n=100$; $n$ "--- число членов ряда (\ref{kud:eq32})

\label{kud:pic11}}

\\

&\vspace{-3.4cm}

[Figure~\ref{kud:pic11}. The change of the rod displacement over time for various~$\xi$ when $\operatorname{\sf Fo}_ 1^c \hm =\operatorname{\sf Fo}_2^c =0.1$, $\operatorname{\sf Fo}_ 3^c =0$, $n=100$; 

$n$ is the number \centerline{of terms in the series (\ref{kud:eq32})]}



\end{tabular}



\vspace{-5mm}



\end{figure}









На рис.~\ref{kud:pic11} приведены результаты расчетов перемещений по формуле \eqref{kud:eq32} при

$\operatorname{\sf Fo}_ 1^c =\operatorname{\sf Fo}_2^c =10$, $\operatorname{\sf Fo}_ 3^c =0$. 

Их анализ позволяет заключить, что с увеличением $\operatorname{\sf Fo}_ 1 ^c$ и~$\operatorname{\sf Fo}_2^c

$ (при $\operatorname{\sf Fo}_3^c =0$) амплитуда колебаний уменьшается. 

При этом в~каждой точке стержня до момента его возвращения в исходное

(недеформированное) состояние совершаются многократные колебательные

процессы в области лишь положительных значений перемещений. 

С~увеличением времени ($\operatorname{\sf Fo}\ge 35$) процесс колебаний стабилизируется вблизи нулевого значения перемещения.



Найдем распределение напряжений в стержне по формуле закона Гука без учета ($\tau_1^c = \tau_2^c =0$) и с учетом ($ \tau_1^c \ne 0$, $ \tau_2^c \ne 0$) релаксационных слагаемых. 

Закон Гука (соотношение \eqref{kud:eq11} при $ \tau_1 ^c= \tau_2 ^c=0$) с учетом

обозначений $ \Theta\hm=U/U_0 $ и $\xi=x/\delta$ будет иметь следующий вид:

\begin{equation}

\label{kud:eq33}

\bar \sigma ( \xi,\operatorname{\sf Fo}) =\frac{\partial  \Theta ( \xi, \operatorname{\sf Fo})} {\partial  \xi},

\end{equation}

где $$ \bar \sigma (\xi,\operatorname{\sf Fo})= \frac {\sigma (x, t)}E \frac \delta{U_0},$$ $\Theta( \xi, \operatorname{\sf Fo})$ "--- точное аналитическое решение уравнения незатухающих колебаний:

\begin{equation}

\label{kud:eq34}

\Theta( \xi,\operatorname{\sf Fo})= \frac{8}{\pi^2} 

\sum\limits_{k=1}^\infty \frac{1}{(2k-1)^2} \cos \Bigl(\frac{2k-1}2 \pi \operatorname{\sf Fo} \Bigr)

\cos \Bigl(\frac{2k-1}2 \pi \xi \Bigr).

\end{equation}



Результаты расчетов безразмерного напряжения по формуле \eqref{kud:eq33} с учетом формулы \eqref{kud:eq34} приведены на рис.~\ref{kud:pic12}. 

Их анализ позволяет заключить, что колебания $\bar \sigma$ происходят скачкообразно во времени от $\bar \sigma=- 0.1$ до $ \bar \sigma=0.1$.



Выпишем соотношение \eqref{kud:eq11} в случае, когда $\tau_1^c \ne 0$ и $ \tau_2^c \ne 0$, в безразмерном виде:

\begin{equation}

\label{kud:eq35}

\frac{ \partial \bar \sigma ( \xi,\operatorname{\sf Fo}) }{\partial \operatorname{\sf Fo}}+A_1  {\bar \sigma}( \xi,\operatorname{\sf Fo}) =A_2 ( \xi,\operatorname{\sf Fo}),

\end{equation}

где 

$$

A_1 = \frac{1}{\operatorname{\sf Fo}_ 1^c}, 

\quad A_2

( \xi,\operatorname{\sf Fo})=\frac{ 1}{\operatorname{\sf Fo}_ 1^c }\frac{\partial  \Theta( \xi,\operatorname{\sf Fo})}{\partial  \xi}+\frac{\operatorname{\sf Fo}_ {2}^c}{\operatorname{\sf Fo}_ 1^c}\frac{\partial

^2 \Theta( \xi,\operatorname{\sf Fo})}{\partial

 \xi\partial \operatorname{\sf Fo}}.

$$





\begin{figure}[b!]



\footnotesize \centering



\begin{tabular}{p{7.7cm}p{4.8cm}}

\includegraphics[scale=0.8]{kudinov12}~~~&

\vspace{-5.25cm}

\caption{(онлайн в цвете) Изменение напряжений в стержне во времени при различных значениях $\xi$, когда 

$\operatorname{\sf Fo}_ 1^c =\operatorname{\sf Fo}_2^c =0$ (расчет по формуле (\ref{kud:eq33}))

\label{kud:pic12}}

\\

&\vspace{-3.4cm}

[Figure~\ref{kud:pic12}. (Color online) The stress change in the rod depending on  dimensionless time ({\sf Fo}) for various $\xi$ when $\operatorname{\sf Fo}_ 1^c =\operatorname{\sf Fo}_2^c =0$  \centerline{(calculation by  Eq.~(\ref{kud:eq33}))]}

\\



\includegraphics[scale=0.8]{kudinov13}~~~&

\vspace{-5.25cm}

\caption{(онлайн в цвете) Изменение напряжений в стержне во времени при различных значениях $\xi$, когда 

$\operatorname{\sf Fo}_ 1^c =\operatorname{\sf Fo}_2^c =0.1$

(расчет по формуле (\ref{kud:eq36}))

\label{kud:pic13}}

\\

&\vspace{-3.4cm}

[Figure~\ref{kud:pic13}. (Color online) The stress change in the rod depending on  dimensionless time ({\sf Fo}) for various $\xi$ when $\operatorname{\sf Fo}_ 1^c =\operatorname{\sf Fo}_2^c =0.1$ 

\centerline{(calculation by  Eq.~(\ref{kud:eq36}))]}

\\





\end{tabular}



\vspace{-5mm}



\end{figure}





Интегрируя уравнение \eqref{kud:eq35}, находим

\begin{equation}

\label{kud:eq36}

 {\bar \sigma}( \xi,\operatorname{\sf Fo})=

 \exp (-A_1 \operatorname{\sf Fo}) 

 \biggl( \int A_2 (\xi, \operatorname{\sf Fo})\exp (A_1 \operatorname{\sf Fo})d\operatorname{\sf Fo}+ C \biggr) ,

\end{equation}

где $C$ "--- постоянная интегрирования, определяемая из начального условия

$\bar \sigma(\xi, 0)= 0$.

В качестве функции $ \Theta( \xi,\operatorname{\sf Fo})$ в формуле \eqref{kud:eq36} использовалось решение \eqref{kud:eq32}.



Анализ результатов расчетов напряжения по формуле \eqref{kud:eq36} позволяет заключить,

что напряжения различных точек стержня изменяются во времени по

гармонической зависимости с затухающей амплитудой (рис.~\ref{kud:pic13}). Таким образом,

учет релаксационных слагаемых позволяет устранить скачкообразное изменение

напряжений во времени.







\smallskip



\Section{Резонансные колебания стержня с учетом релаксационных свойств

материалов}

Если в задаче \eqref{kud:eq29}--\eqref{kud:eq31} к свободному торцу стержня приложить нагрузку,

изменяющуюся по гармоническому закону, то первое граничное условие из \eqref{kud:eq31}

примет вид

\begin{equation}

\label{kud:eq37}

\frac{\partial  \Theta(\xi ,\operatorname{\sf Fo})} {\partial  \xi}\Big|_{\xi=0}=\operatorname{\sf Fo}_4 \cos

(\operatorname{\sf Fo}_5 \operatorname{\sf Fo}), \vspace{-2mm}

\end{equation}

где $\operatorname{\sf Fo}_4 = \delta /U_0$; $\operatorname{\sf Fo}_5 = \omega \delta

/e_s$; $ \omega =2 \pi \nu_l $ "--- круговая частота;

$ \nu_l$ "--- частота колебаний нагрузки.



Далее будет показано, что при совпадении частоты собственных колебаний

стержня с частотой $\operatorname{\sf Fo}_5 $ внешней нагрузки наблюдаются резонансные

колебания стержня, в том числе и бифуркационные (флаттерные).



Получение точного аналитического решения краевой задачи \eqref{kud:eq29}--\eqref{kud:eq31} \linebreak с~учетом граничного условия \eqref{kud:eq37} существенно усложняется, и для ее решения был

использован конечно"=разностный метод.



Результаты выполненных исследований позволяют сделать вывод, что при

$\operatorname{\sf Fo}_ 1^c =\operatorname{\sf Fo}_ {2}^c =\operatorname{\sf Fo}_3^c =0$ колебания

незатухающие. 

При $\operatorname{\sf Fo}_3^c =0.1$ и $\operatorname{\sf Fo}_ 1^c \hm

=\operatorname{\sf Fo}_ {2}^c =\operatorname{\sf Fo}_4^c =0$ колебания становятся затухающими с~экспоненциально уменьшающейся амплитудой. 

При $\operatorname{\sf Fo}_ 1^c \hm =\operatorname{\sf Fo}_ {2}^c =\operatorname{\sf Fo}_3 ^c =0$ и 

$\operatorname{\sf Fo}_4 =0.5$, $\operatorname{\sf Fo}_5 \hm =1.575$ безразмерная частота собственных колебаний стержня $\nu_l =1.575$ совпадает с частотой вынужденных колебаний, возникающих под действием внешней нагрузки вида \eqref{kud:eq37}. 

При этом наблюдается неограниченное возрастание амплитуды колебаний (рис.~\ref{kud:pic14}).



С увеличением коэффициента сопротивления среды $\operatorname{\sf Fo}_3^c $ амплитуда

колебаний внутренних точек стержня уменьшается (рис.~\ref{kud:pic15}) и при некотором

большом его значении ($\operatorname{\sf Fo}_3^c \ge 100)$ возврат стержня в исходное

состояние происходит практически при отсутствии колебательного процесса внутренних точек стержня (критическое затухание) при неизменной амплитуде колебаний внешней нагрузки.



Результаты расчетов резонансных колебаний для $\operatorname{\sf Fo}_ 1^c =\operatorname{\sf Fo}_2^c =2$, $\operatorname{\sf Fo}_3^c =1$, $\operatorname{\sf Fo}_4 =0.1$, $\operatorname{\sf Fo}_5 \hm

=1.575$ приведены на рис.~\ref{kud:pic16}. 

При дальнейшем увеличении коэффициентов релаксации ($\operatorname{\sf Fo}_ 1^c =\operatorname{\sf Fo}_2^c =10$, $\operatorname{\sf Fo}_3^c =0.3$, $\operatorname{\sf Fo}_4 =0.1$, $\operatorname{\sf Fo}_5 =1.575$) в процессе колебаний наблюдается бифуркационный резонанс, при котором амплитуда колебаний возрастает периодически. 

При этом в зависимости от соотношения величин $\operatorname{\sf Fo}_ 1^c$, $\operatorname{\sf Fo}_ {2} ^c$ и $\operatorname{\sf Fo}_3^c $ процесс колебаний может быть как затухающим, так и незатухающим.



Отметим, что бифуркационное изменение амплитуды колебаний наблюдается также

на частотах колебаний внешней нагрузки, близких к частоте собственных

колебаний стержня, но не совпадающих с ней полностью. Результаты расчетов

для такого типа колебаний приведены на рис.~\ref{kud:pic17}, \ref{kud:pic18}.



Обнаруженные в результате исследований бифуркационные (флаттерные) колебания

как в резонансных, так и (что особенно важно) в нерезонансных частотах

представляют серьезную проблему при работе различных технических устройств и

приводят к значительному сокращению их ресурса. Результаты исследований

позволяют найти такое соотношение частот собственных и вынужденных

колебаний, при которых бифуркационные колебания могут быть исключены, то

есть эти результаты позволят разработать рекомендации по установке

демпфирующих устройств, способствующих затуханию колебаний, минуя

бифуркационный режим.







\begin{figure}[p!]



\centering \footnotesize



\begin{tabular}{rp{4.8cm}}

\includegraphics[scale=0.8]{kudinov14}\!\!&

\vspace{-5.25cm}

\caption{(онлайн в цвете) Изменение перемещений стержня в резонансных частотах при различных значениях $\xi$, когда  $\operatorname{\sf Fo}_ 1^c =\operatorname{\sf Fo}_ 2^c =\operatorname{\sf Fo}_3^c =0$, 

$\operatorname{\sf Fo}_4 =0.5$, $\operatorname{\sf Fo}_5 = 1.575$

\label{kud:pic14}}

\\

&\vspace{-3.2cm}

[Figure~\ref{kud:pic14}. (Color online)  The change of the rod displacement at the resonant frequencies for various $\xi$ when $\operatorname{\sf Fo}_ 1^c =\operatorname{\sf Fo}_ 2^c ={\operatorname{\sf Fo}_3^c =0}$, 

${\operatorname{\sf Fo}_4 =0.5}$, $\operatorname{\sf Fo}_5 = 1.575$]

\\





\includegraphics[scale=0.8]{kudinov15}\!\!&

\vspace{-5.25cm}

\caption{(онлайн в цвете) Изменение перемещений стержня в резонансных частотах при различных значениях $\xi$, когда  $\operatorname{\sf Fo}_ 1^c =\operatorname{\sf Fo}_ 2 =0$, $\operatorname{\sf Fo}_3 ^c =100$, ${\operatorname{\sf Fo}_4 =1}$, $\operatorname{\sf Fo}_5 =1.575$

\label{kud:pic15}}

\\

&\vspace{-3.2cm}

[Figure~\ref{kud:pic15}. (Color online)  The change of the rod displacement at the resonant frequencies for various $\xi$ when $\operatorname{\sf Fo}_ 1^c =\operatorname{\sf Fo}_ 2 =0$, $\operatorname{\sf Fo}_3 ^c =100$, $\operatorname{\sf Fo}_4 =1$, $\operatorname{\sf Fo}_5 =1.575$]

\\



\includegraphics[scale=0.8]{kudinov16}\!\!\!\!&

\vspace{-5.25cm}

\caption{(онлайн в цвете) Изменение перемещений стержня в резонансных частотах при различных значениях $\xi$, когда  $\operatorname{\sf Fo}_ 1^c =\operatorname{\sf Fo}_2^c =2$, $\operatorname{\sf Fo}_3 ^c =1$, ${\operatorname{\sf Fo}_4 =0.1}$, $\operatorname{\sf Fo}_5 =1.575$

\label{kud:pic16}}

\\

&\vspace{-3.2cm}

[Figure~\ref{kud:pic16}. (Color online)  The change of the rod displacement at the resonant frequencies for various $\xi$ when $\operatorname{\sf Fo}_ 1^c =\operatorname{\sf Fo}_2^c =2$, $\operatorname{\sf Fo}_3 ^c =1$, $\operatorname{\sf Fo}_4 =0.1$, $\operatorname{\sf Fo}_5 =1.575$]

\\



\end{tabular}





\end{figure}







\begin{figure}[t!]



\centering \footnotesize



\begin{tabular}{rp{4.8cm}}

\includegraphics[scale=0.8]{kudinov17}\!\!&

\vspace{-5.25cm}

\caption{(онлайн в цвете) Изменение перемещений стержня в нерезонансных частотах при различных значениях $\xi$, когда  $\operatorname{\sf Fo}_ 1^c =\operatorname{\sf Fo}_2^c =10$, $\operatorname{\sf Fo}_3 ^c=0.3$, $\operatorname{\sf Fo}_4 =0.1$, $\operatorname{\sf Fo}_5 =0.1$

\label{kud:pic17}}

\\

&\vspace{-3.2cm}

[Figure~\ref{kud:pic17}. (Color online)  The change of the rod displacement at the nonresonant frequencies for various $\xi$ when $\operatorname{\sf Fo}_ 1^c =\operatorname{\sf Fo}_2^c =10$, $\operatorname{\sf Fo}_3 ^c=0.3$, $\operatorname{\sf Fo}_4 =0.1$, $\operatorname{\sf Fo}_5 =0.1$]

\\





\includegraphics[scale=0.8]{kudinov18}\!\!&

\vspace{-5.25cm}

\caption{(онлайн в цвете) Изменение перемещений стержня в нерезонансных частотах при различных значениях $\xi$, когда  

$\operatorname{\sf Fo}_ 1^c =\operatorname{\sf Fo}_2 =10$,  

$\operatorname{\sf Fo}_3 ^c=0.3$, $\operatorname{\sf Fo}_4 =0.1$, $\operatorname{\sf Fo}_5 =1.5$

\label{kud:pic18}}

\\

&\vspace{-3.2cm}

[Figure~\ref{kud:pic18}. (Color online)  The change of the rod displacement at the nonresonant frequencies for various $\xi$ when $\operatorname{\sf Fo}_ 1^c =\operatorname{\sf Fo}_2 =10$,  

$\operatorname{\sf Fo}_3 ^c=0.3$, $\operatorname{\sf Fo}_4 =0.1$, $\operatorname{\sf Fo}_5 =1.5$]

\\





\end{tabular}





\end{figure}











\Section{Математическая модель динамической термоупругости с учетом

пространственно"=временной нелокальности}

При больших скоростях изменения температуры температурные напряжения оказываются зависимыми от распределения напряжений в~предшествующие моменты времени. 

Время в~данном случае становится независимой переменной и~поэтому в~уравнениях термоупругости появляются инерционные члены, учитывающие изменение напряжений во времени~[\citen{kud:bib18}--\citen{kud:bib21}].



Классическое одномерное уравнение динамической термоупругости для незакрепленной пластины имеет вид~[\citen{kud:bib18}--\citen{kud:bib20}]

\begin{equation}

\label{kud:eq38}

\frac{\partial^2\sigma_x (x,t)}{\partial x^2}-\frac{1}{c_w^2}\frac{\partial^2\sigma _x (x,t)}{\partial t^2}=\frac{1+\nu }{1-\nu}\rho \alpha \frac{\partial ^2T(x,t)}{\partial t^2}, \quad

t>0, \quad 0<x< \delta,

\end{equation}

где $c_w^2=E(1- \nu)/[ (1+ \nu)(1-2 \nu )\rho ]$, 

$c_w$ "--- скорость распространения упругой волны, 

$\delta $ "--- толщина пластины, 

$\sigma _x$ "--- нормальное напряжение в направлении координаты

$x$, $T$ "--- температура, 

$T_0$ "--- температура в ненапряженном состоянии, 

$t$ "--- время, $E$ "--- модуль упругости, $\nu$ "--- коэффициент Пуассона, 

$\alpha$ "--- коэффициент линейного расширения, $\rho$ "--- плотность.%, $U$ "--- перемещение, 

%$ \varepsilon_x \hm =\partial U/\partial x$ "--- деформация.



Распределение температуры как функции пространственной переменной и~времени предполагается известным из решения соответствующей задачи теплопроводности.



Уравнение \eqref{kud:eq37} является волновым, описывающим незатухающие колебания температурных напряжений. Отсутствие затухания связано с~тем, что данное уравнение не содержит членов, учитывающих силы внутреннего сопротивления среды процессу изменения ее формы в результате температурных деформаций. 

Для учета сопротивления среды будем использовать уравнение равновесия (движения) вида

\begin{equation}

\label{kud:eq39}

\frac{\partial \sigma_x }{\partial x}-\gamma \rho \frac{\partial U}{\partial t}= {\rho }\frac{\partial^2U}{\partial

t^2},

\end{equation}

где второе слагаемое в левой части представляет силу сопротивления, приходящуюся на единицу объема среды,  $\gamma$ "--- коэффициент сопротивления.%, $U$ "--- перемещение.





Уравнение \eqref{kud:eq37} выведено без учета пространственно"=временной локальной

неравновесности. Для ее учета при выводе уравнения динамической

термоупругости будем использовать формулу закона Гука вида \eqref{kud:eq11}. Подставляя

\eqref{kud:eq11} в \eqref{kud:eq38} с учетом соотношения между напряжением и деформацией, находим

следующее уравнение динамической термоупругости (в безразмерном виде):

\begin{gather}

\frac{\partial^2 \bar \sigma( \xi,\operatorname{\sf Fo})}{\partial  \xi^2}+\operatorname{\sf Fo}_1^p

\frac{\partial^3 \bar \sigma( \xi,\operatorname{\sf Fo})}{\partial  \xi^2\partial \operatorname{\sf Fo}}-

\frac{\partial^2 \bar \sigma( \xi,\operatorname{\sf Fo})}{\partial \operatorname{\sf Fo}^2}-\operatorname{\sf Fo}_3^p 

\frac{\partial  \bar \sigma( \xi,\operatorname{\sf Fo})}{\partial \operatorname{\sf Fo}}=\\

\label{kud:eq41}

\quad =

\operatorname{\sf Fo}_2^p \frac{\partial^3\Theta ( \xi,\operatorname{\sf Fo})}{\partial  \xi^2\partial \operatorname{\sf Fo}}+

\frac{\partial

^2\Theta ( \xi,\operatorname{\sf Fo})}{\partial

\operatorname{\sf Fo}^2}+\operatorname{\sf Fo}_3^p \frac{\partial \Theta ( \xi,\operatorname{\sf Fo})}{\partial \operatorname{\sf Fo}},\\

\operatorname{\sf Fo}>0, \quad  0< \xi<1, \quad  0\le  \Theta\le 1.

\end{gather}

Здесь

$$

\operatorname{\sf Fo}_1^p =\frac{c}{2\delta}\Bigl(  \tau_2 -\frac{ \tau_1 E}{ c^2 \rho } \Bigr),

\quad 

\operatorname{\sf Fo}_2^p = \frac12 \frac{ \tau_1 E}{\delta c\rho },

\quad 

\operatorname{\sf Fo}_3 ^p = \frac{2 c \rho \gamma \delta}E \frac{(1+ \nu )(1-2 \nu

)}{1- \nu },

$$ 

%безразмерные параметры; $A=$; $B=\frac{(1+ \nu)}{(1-\nu)} \alpha$; 

$$ \Theta=\frac{T-T_{w}}{\Delta T}, \quad 

\xi=\frac{x}{2\delta}, \quad 

\operatorname{\sf Fo}=\frac{ct}{2\delta }, \quad \bar  \sigma= \frac{( 1-2 \nu) \sigma_x }{E \alpha

\Delta T}, \quad 

\Delta T=T_0 -T_{w};$$ 

$\bar \sigma$, $ \Theta$, $ \xi $ "--- соответственно безразмерные напряжение, температура, координата; $\operatorname{\sf Fo}$ "--- число Фурье (безразмерное время); $T_0$ "--- начальная температура; $T_{w}$ "--- температура стенки пластины; $E$ "--- модуль упругости; $\delta$ "--- толщина пластины.



Начальные и граничные условия к уравнению \eqref{kud:eq39} для бесконечной пластины при

отсутствии внешней нагрузки имеют вид

\begin{gather}

\label{kud:eq41}

\sigma(\xi,0)=0,

\quad

\frac{\partial  \sigma( \xi, \operatorname{\sf Fo})}{\partial \operatorname{\sf Fo}}\Bigr|_{\operatorname{\sf Fo}=0}=0;\\

\label{kud:eq41a}

\sigma(0,\operatorname{\sf Fo})=0,

\quad

\sigma(1,\operatorname{\sf Fo})=0.

\end{gather}



Очевидно, что при $\operatorname{\sf Fo}_ 1 ^p =\operatorname{\sf Fo}_2^p

=\operatorname{\sf Fo}_3^p =0$ задача \eqref{kud:eq39}--\eqref{kud:eq41a} приводится к~задаче динамической

термоупругости без учета пространственно"=временной нелокальности и~сопротивления среды.



В качестве решения температурной задачи использовались приведенные выше

точные аналитические решения задачи теплопроводности с~учетом

пространственно"=временной нелокальности.



Точное аналитическое решение задачи \eqref{kud:eq39}--\eqref{kud:eq41a} имеет вид

\begin{multline}

\label{kud:eq42}

\bar \sigma ( \xi,\operatorname{\sf Fo} )=

\sum\limits_{k=1}^\infty \Bigl(\bigl(C_{1k} e^{-\nu_{ 1k} \operatorname{\sf Fo}}+C_{2k} e^{-\nu_{2k} \operatorname{\sf Fo}}\bigr)  -\\

- \frac1{H_3 H_4} \bigl(A_2 H_3  e^{ z_{ 1k} {m}\operatorname{\sf Fo}}+B_2

H_4  e^{z_{2k} { m} \operatorname{\sf Fo}}\bigr)\Bigr)\cos  

\Bigl(\frac{(2k-1)}{2}(1-2 \xi)\pi \Bigr),

\end{multline}

где 

$$

\nu_{1,2\,k} =\frac 12 \Bigl(\operatorname{\sf Fo}_{4k} \pm \sqrt { \operatorname{\sf Fo}_{4k}^2

-4 {\mu }_k } \Bigr),

\quad  

z_{1k} =- \mu_k =(2k-1)^2 \pi^2,

\quad 

z_{2k} =-\frac1{\operatorname{\sf Fo}_2^p};

$$ 

$C_{1k}$, $C_{2k}$ "--- константы интегрирования,

определяемые из начальных условий \eqref{kud:eq41}; $m$, $A_2$, $B_2$, $H_3$, $H_4$ "---

заданные постоянные; $\operatorname{\sf Fo}_{4k}=\mu_k \operatorname{\sf Fo}_ 1^p

\operatorname{\sf Fo}_2^p$.





Результаты расчетов температурных напряжений по формуле \eqref{kud:eq42} приведены на

рис.~\ref{kud:pic19}--\ref{kud:pic22}. 

В~частности, на рис.~\ref{kud:pic19} даны результаты расчетов при 

$\operatorname{\sf Fo}_1^p =\operatorname{\sf Fo}_2^p =\operatorname{\sf Fo}_3^p =0$ 

для значений $\xi=0.1$, 0.3, 0.5. 

Их анализ позволяет заключить, что напряжения в этих

точках с течением времени изменяются скачкообразно, периодически изменяя

знак. 

Волновой процесс изменения напряжений при 

$\operatorname{\sf Fo}_1^p =\operatorname{\sf Fo}_2^p

=\operatorname{\sf Fo}_3^p =0$ является незатухающим во времени, что объясняется

отсутствием в уравнении \eqref{kud:eq39} членов, отвечающих затуханию колебаний.

Отметим, что при $\operatorname{\sf Fo}_3^p \ne 0$ в уравнении \eqref{kud:eq39} появляется слагаемое,

которое учитывает силы внутреннего трения.







\begin{figure}[p!]



\centering \footnotesize



\begin{tabular}{lp{4.8cm}}

\includegraphics[scale=0.8]{kudinov19}~\,&

\vspace{-5.25cm}

\caption{(онлайн в цвете) Изменение напряжений в пластине для различных $\xi$ (расчет по формуле (\ref{kud:eq42})), когда $\operatorname{\sf Fo}_1^p\hm =\operatorname{\sf Fo}_2^p =\operatorname{\sf Fo}_3^p =0$, $m=0.001$, 

$\operatorname{\sf Fo}_r=0.001$ "--- коэффициент релаксации из решения температурной задачи; $n=1000$  "--- число членов ряда

\label{kud:pic19}}

%\\

%&%\vspace{-2.8cm}

%[Figure~\ref{kud:pic19}. (Color online)  The change of stresses in the plate for

%various $\xi$ (calculation by Eq.~(\ref{kud:eq42})) when ${\operatorname{\sf Fo}_1^p =\operatorname{\sf Fo}_2^p =\operatorname{\sf Fo}_3^p =0}$, $m=0.001$, 

%$\operatorname{\sf Fo}_r=0.001$, $n=1000$; $\operatorname{\sf Fo}_r$ is the relaxation coefficient from the solution of the temperature problem; $n$ is the number of terms in the 

%\centerline{series~(\ref{kud:eq42})]}

\\

\multicolumn{2}{p{\textwidth}}{

[Figure~\ref{kud:pic19}. (Color online)  The change of stresses in the plate for

various $\xi$ (calculation by Eq.~(\ref{kud:eq42})) when ${\operatorname{\sf Fo}_1^p =\operatorname{\sf Fo}_2^p =\operatorname{\sf Fo}_3^p =0}$, $m=0.001$, 

$\operatorname{\sf Fo}_r=0.001$, $n=1000$; $\operatorname{\sf Fo}_r$ is the relaxation coefficient from the solution of the temperature problem; $n$ is the number of terms in the 

\centerline{series~(\ref{kud:eq42})]}}

\\



\\



\includegraphics[scale=0.8]{kudinov20}~&

\vspace{-5.25cm}

\caption{(онлайн в цвете) Распределение напряжений в пластине для различных чисел $\sf Fo$ (расчет по формуле (\ref{kud:eq42})), когда $\operatorname{\sf Fo}_1^p =\operatorname{\sf Fo}_2^p =\operatorname{\sf Fo}_3^p =0$, $m=0.001$, 

$\operatorname{\sf Fo}_r=0.001$ "--- коэффициент релаксации из решения температурной задачи; $n=1000$  "--- число членов ряда

\label{kud:pic20}}

\\ 

\multicolumn{2}{p{\textwidth}}{

[Figure~\ref{kud:pic20}. (Color online)  The stress distribution in the plate for different Fourier numbers

($\sf Fo$) (calculation by Eq.~(\ref{kud:eq42})) when $\operatorname{\sf Fo}_1^p =\operatorname{\sf Fo}_2^p =\operatorname{\sf Fo}_3^p =0$, $m=0.001$, 

$\operatorname{\sf Fo}_r=0.001$, $n=1000$; $\operatorname{\sf Fo}_r$ is the relaxation coefficient from the solution of the temperature problem; $n$ is the number of \centerline{terms in the 

series~(\ref{kud:eq42})]}}

\\



\\

\includegraphics[scale=0.8]{kudinov21}~&

\vspace{-6.35cm}

\caption{(онлайн в цвете) Распределение напряжений в точке $\xi=0.5$ для различных $\operatorname{\sf Fo}_1^p =\operatorname{\sf Fo}_2 ^p$ (расчет по формуле (\ref{kud:eq42})), когда 

$\operatorname{\sf Fo}_3^p =0.1$, $m=0.001$, 

$\operatorname{\sf Fo}_r \hm =0.001$ "--- коэффициент релаксации из решения температурной задачи;  $n=1000$  "--- число членов ряда

\label{kud:pic21}}

\\ 

&\vspace{-2.9cm}

[Figure~\ref{kud:pic21}. (Color online)  The stress distribution at the point  $\xi=0.5$ for different $\operatorname{\sf Fo}_1^p =\operatorname{\sf Fo}_2 ^p$ (calculation by Eq.~(\ref{kud:eq42})) when $\operatorname{\sf Fo}_3^p =0.1$, $m=0.001$, 

$\operatorname{\sf Fo}_r=0.001$, $n=1000$; $\operatorname{\sf Fo}_r$ is the relaxation coefficient from the solution of the temperature problem;  $n$ is the number of terms \centerline{in the series (\ref{kud:eq42})]}

\\



\end{tabular}



\end{figure}









\begin{figure}[t!]



\centering \footnotesize



\begin{tabular}{lp{4.8cm}}

\includegraphics[scale=0.8]{kudinov22}~\,&

\vspace{-5.25cm}

\caption{(онлайн в цвете) Распределение напряжений в пластине для различных чисел $\sf Fo$ (расчет по формуле (\ref{kud:eq42})), когда $\operatorname{\sf Fo}_1^p\hm =\operatorname{\sf Fo}_2^p =\operatorname{\sf Fo}_3^p \hm =0.01$, $m\hm =0.001$, 

$\operatorname{\sf Fo}_r=0.001$ "--- коэффициент релаксации из решения температурной задачи;  ${n=1000}$  "--- число членов ряда

\label{kud:pic22}}

\\

%&\vspace{-2.8cm}

%[Figure~\ref{kud:pic22}. (Color online)  The stress distribution in the plate for different Fourier numbers

%($\sf Fo$) (calculation by Eq.~(\ref{kud:eq42})) when $\operatorname{\sf Fo}_1^p =\operatorname{\sf Fo}_2^p =\operatorname{\sf Fo}_3^p =0.01$, $m=0.001$, 

%$ \operatorname{\sf Fo}_r=0.001$, $n=1000$; $n$ is the number of terms in the 

%series~(\ref{kud:eq42})]

\multicolumn{2}{p{\textwidth}}{

[Figure~\ref{kud:pic22}. (Color online)  The stress distribution in the plate for different Fourier numbers

($\sf Fo$) (calculation by Eq.~(\ref{kud:eq42})) when $\operatorname{\sf Fo}_1^p =\operatorname{\sf Fo}_2^p =\operatorname{\sf Fo}_3^p =0.01$, $m=0.001$, 

$ \operatorname{\sf Fo}_r=0.001$, $n=1000$; $\operatorname{\sf Fo}_r$ is the relaxation coefficient from the solution of the temperature problem;   $n$ is the number of \centerline{terms in the 

series~(\ref{kud:eq42})]}

}

\\



\end{tabular}



\vspace{-5mm}



\end{figure}







На рис.~\ref{kud:pic20} приведены результаты расчетов распределения напряжений при

$\operatorname{\sf Fo}_1^p =\operatorname{\sf Fo}_2^p =\operatorname{\sf Fo}_3^p =0$ по координате $\xi$ для

различных моментов времени. 

Их анализ позволяет заключить, что изменение напряжений по форме эквивалентно случаю колебаний закрепленной с~двух сторон струны, имеющей изломы (в виде прямоугольных скачков напряжений),

перемещающиеся по координате $\xi$ во времени. 

Отметим, что результаты расчетов, приведенные на рис.~\ref{kud:pic19}, \ref{kud:pic20}, совпадают с~результатами,

полученными в~работе~\cite{kud:bib21}.



Результаты расчетов напряжений в точке $\xi=0.5$ для различных значений $\operatorname{\sf Fo}_1^p =\operatorname{\sf Fo}_2^p $  и~при $\operatorname{\sf Fo}_3^p =0.1$ приведены на рис.~\ref{kud:pic21}. 

Их анализ позволяет заключить, что процесс колебаний является затухающим во времени. 

При увеличении значений $\operatorname{\sf Fo}_1^p =\operatorname{\sf Fo}_2^p$ (при неизменном $\operatorname{\sf Fo}_3^p =0.1)$ скачки напряжений сглаживаются и~при некоторых больших значениях $\operatorname{\sf Fo}_1^p =\operatorname{\sf Fo}_2^p$ возникающие от теплового удара напряжения монотонно уменьшаются до нулевого значения при практическом отсутствии колебательного процесса.



На рис.~\ref{kud:pic22} приведены результаты расчетов напряжений по координате $\xi$ для $0.1\le \operatorname{\sf Fo} \le 1.8$. 

Их анализ позволяет заключить, что колебания напряжений по форме аналогичны колебаниям закрепленной на концах струны, не имеющей изломов.







Если в качестве решения температурной задачи использовать точное

аналитическое решение классического параболического уравнения

теплопроводности, то решение динамической задачи термоупругости приводит к

бесконечным температурным напряжениям в точках задания граничного условия

первого рода (тепловой удар), что связано со следующими обстоятельствами.

При выполнении условий теплового удара предполагается мгновенное во времени

принятие поверхностью пластины температуры, задаваемой граничным условием

первого рода. Так как в динамической задаче термоупругости учитывается

изменение термических напряжений в зависимости от скорости изменения

температуры во времени, то мгновенное ее задание означает бесконечную

скорость изменения температуры, что и приводит к бесконечным значениям

температурных напряжений.



Если использовать в качестве решения температурной задачи точное

аналитическое решение краевой задачи теплопроводности с учетом

пространственно"=временномй нелокальности, то напряжения в точке задания

граничных условий первого рода имеют конечную величину. Этот факт можно

объяснить тем, что температура стенки в данном случае не может быть задана

мгновенно, а только лишь в течение некоторого конечного промежутка времени.



\smallskip 



\Section{Исследование локально"=неравновесных процессов теплового вос\-пла\-ме\-нения}

В теории теплового воспламенения изучается протекание химической реакции в

условиях прогрессивного самоускорения, приводящего к неограниченному

возрастанию температуры (тепловому воспламенению)~\cite{kud:bib22,kud:bib23}. В~уравнениях

известных математических моделей теплового воспламенения не учитываются

релаксационные свойства материалов (локальная неравновесность процесса),

что, как показали приведенные ниже исследования, приводит к существенному

отклонению критических условий и времени задержки теплового воспламенения по

сравнению со случаем их учета. Данный факт можно объяснить тем, что в основе

известных моделей лежат параболические уравнения, выведенные из условия

бесконечной скорости распространения теплоты. В работах~[\citen{kud:bib9}--\citen{kud:bib12}] приводятся

результаты построения моделей, в которых учитывается ускорение во времени

как теплового потока, так и градиента температуры. В настоящей работе

теоретические положения построения таких моделей были распространены

применительно к решению краевой задачи теплового воспламенения для

бесконечной пластины с экспоненциально зависящим от температуры внутренним

источником теплоты 

$$\omega (T)=Q\rho k_0 \exp \Bigl(- \frac{E}{RT} \Bigr)$$ в следующей

математической постановке:

\begin{gather}

\frac{\partial \Theta }{\partial \operatorname{\sf Fo}}+

\operatorname{\sf Fo}_ 1 \frac{\partial^2\Theta }{\partial \operatorname{\sf Fo}^2}=

\frac{\partial ^2\Theta }{\partial \xi ^2}+\operatorname{\sf Fo}_ 2 

\frac{\partial ^3\Theta }{\partial \xi ^2\partial \operatorname{\sf Fo}}-

\frac{\gamma }{\left(\Theta \Delta T+T_{w} \right)^2}\times

\\

\label{kud:eq43}

\times \exp\Bigl( -\frac{E}{R(\Theta \Delta T+T_{w})} \Bigr)

\frac{\partial \Theta}{\partial \operatorname{\sf Fo}}

+

\tau^p \exp\Bigl( -\frac{E}{R(\Theta \Delta T+T_{w})}\Bigr), \\ 

\operatorname{\sf Fo}>0, \quad 0<\xi <1;

\\

\label{kud:eq44}

\Theta (\xi,0)=1,

\;\;

\partial 

\frac{\Theta (\xi,\operatorname{\sf Fo})}{\partial \operatorname{\sf Fo}}\Big|_{\operatorname{\sf Fo}=0} =0;

\quad

\frac{\partial \Theta (\xi,\operatorname{\sf Fo})}{\partial \xi}\Big|_{\xi=0} =0,

\; \;

\Theta (1,\operatorname{\sf Fo})=0,

\end{gather}

где 

$$

\Theta =\frac{T-T_{w}}{\Delta T},

\quad  \operatorname{\sf Fo}=\frac{at}{\delta^2}, 

\quad \xi =\frac{x}{\delta },

\quad \operatorname{\sf Fo}_1 =\frac{a\tau _1}{\delta^2},

\quad \operatorname{\sf Fo}_2 =\frac{a\tau _2 }{\delta ^2},

$$ 

$$

\tau^p =\frac{\delta^2Qk_0 }{ca\Delta T}, 

\quad \Delta T=T_0 -T_{w},

\quad \gamma = \frac{\tau _1 k_0 E}{cR};

$$

 $\Theta$, $\operatorname{\sf Fo}$, $\xi$ "--- соответственно безразмерные температура,

время, координата; 

$\tau^p $ "--- безразмерный параметр; 

$\operatorname{\sf Fo}_1$, $\operatorname{\sf Fo}_2 $ "--- безразмерные коэффициенты релаксации; $T$ "---

температура; $x$ "--- координата; $a$ "--- коэффициент температуропроводности;

$t$ "--- время; $\omega (T)$ "--- мощность внутреннего источника теплоты; 

$\rho $ "--- плотность; 

$c$ "--- теплоемкость; $Q$ "--- тепловой эффект реакции; $k_0$

"--- предэкспоненциальный множитель; 

$T_0 $ "--- начальная температура;

$T_{w} $ "--- температура стенки; $\tau _1$, $\tau_2$ "---

коэффициенты релаксации теплового потока и градиента температуры в формуле

закона Фурье; $E$ "--- энергия активации; $R$ "--- универсальная газовая

постоянная.



Вывод уравнения \eqref{kud:eq43} базируется на формуле закона Фурье \eqref{kud:eq9}, в которой

учитывается ускорение во времени теплового потока и градиента температуры, и

уравнении теплового баланса

\begin{equation} 

\label{kud:eq45}

c\rho \frac{ \partial T}{\partial t} + \frac{\partial q}{\partial x}+\omega (T)=0,

\end{equation}

где $\lambda $ "--- коэффициент теплопроводности.



Для решения задачи \eqref{kud:eq43}, \eqref{kud:eq44} применялся метод конечных разностей \cite{kud:mc:bib6} со следующими исходными данными: $\tau^p =4.06 \cdot 10^8$; 

$E=127$~кДж/моль; $R=8.31$~Дж/(моль${}\cdot{}$К);

$\Delta T=10$~К; $T_{w} =850$~К; $k_0 =1.925\cdot 10^6$~Гц.



Результаты расчетов приведены на рис.~\ref{kud:pic23}--\ref{kud:pic25}. 

Их анализ позволяет заключить, что величина безразмерных коэффициентов релаксации существенно

влияет на время задержки теплового воспламенения. 

Так, при $\operatorname{\sf Fo}_ 1 =\operatorname{\sf Fo}_2 \hm =0.15$ оно возрастает примерно на~$1/3$ по

сравнению со случаем, когда релаксационные свойства материала не учитываются ($\operatorname{\sf Fo}_ 1 =\operatorname{\sf Fo}_2 =0)$ (см. рис.~\ref{kud:pic23}).



\begin{figure}[p!]





\centering \footnotesize



\begin{tabular}{lp{4.8cm}}

\includegraphics[scale=0.8]{kudinov23}~\,&

\vspace{-5.25cm}

\caption{(онлайн в цвете) Изменение температуры $\Theta$, когда $\xi = 0.15$, $\gamma =0$

\label{kud:pic23}}

\\

&\vspace{-3.8cm} \centering

[Figure~\ref{kud:pic23}. (Color online)  

 The~change of the temperature ($\Theta$) when $\xi = 0.15$, $\gamma =0$]

%\multicolumn{2}{p{\textwidth}}{

%[Figure~\ref{kud:pic23}. (Color online)  The stress distribution in the plate for different Fourier numbers

%($\sf Fo$) (calculation by Eq.~(\ref{kud:eq42})) when $\operatorname{\sf Fo}_1^p =\operatorname{\sf Fo}_2^p =\operatorname{\sf Fo}_3^p =0.01$, $m=0.001$, 

%$ \operatorname{\sf Fo}_r=0.001$, $n=1000$; $\operatorname{\sf Fo}_r$ is the relaxation coefficient from the solution of the temperature problem;   $n$ is the number of \centerline{terms in the 

%series~(\ref{kud:eq42})]}

%}

\\

\end{tabular}



\begin{tabular}{lp{4.8cm}}

\includegraphics[scale=0.8]{kudinov24}~\,&

\vspace{-5.25cm}

\caption{Распределение температуры $\Theta$ для различных чисел $\sf Fo$, когда $\operatorname{\sf Fo}_1 =\operatorname{\sf Fo}_2 =0.15$, $\gamma =0$

\label{kud:pic24}}

\\

&\vspace{-3.8cm} 

[Figure~\ref{kud:pic24}. The temperature distribution for different Fourier numbers ({\sf Fo})  when \centerline{$\operatorname{\sf Fo}_1 =\operatorname{\sf Fo}_2 =0.15$, $\gamma =0$]}

%\multicolumn{2}{p{\textwidth}}{

%[Figure~\ref{kud:pic23}. (Color online)  The stress distribution in the plate for different Fourier numbers

%($\sf Fo$) (calculation by Eq.~(\ref{kud:eq42})) when $\operatorname{\sf Fo}_1^p =\operatorname{\sf Fo}_2^p =\operatorname{\sf Fo}_3^p =0.01$, $m=0.001$, 

%$ \operatorname{\sf Fo}_r=0.001$, $n=1000$; $\operatorname{\sf Fo}_r$ is the relaxation coefficient from the solution of the temperature problem;   $n$ is the number of \centerline{terms in the 

%series~(\ref{kud:eq42})]}

%}

\\



\end{tabular}







\begin{tabular}{lp{4.8cm}}

\includegraphics[scale=0.8]{kudinov25}~\,&

\vspace{-5.25cm}

\caption{(онлайн в цвете) Изменение температуры $\Theta$, когда $\xi = 0.9$

\label{kud:pic25}}

\\

&\vspace{-3.8cm} \centering

[Figure~\ref{kud:pic25}. (Color online)  

 The~change of the temperature ($\Theta$) when $\xi = 0.9$]

%\multicolumn{2}{p{\textwidth}}{

%[Figure~\ref{kud:pic23}. (Color online)  The stress distribution in the plate for different Fourier numbers

%($\sf Fo$) (calculation by Eq.~(\ref{kud:eq42})) when $\operatorname{\sf Fo}_1^p =\operatorname{\sf Fo}_2^p =\operatorname{\sf Fo}_3^p =0.01$, $m=0.001$, 

%$ \operatorname{\sf Fo}_r=0.001$, $n=1000$; $\operatorname{\sf Fo}_r$ is the relaxation coefficient from the solution of the temperature problem;   $n$ is the number of \centerline{terms in the 

%series~(\ref{kud:eq42})]}

%}

\\

\end{tabular}



\end{figure}



Факт задержки теплового воспламенения объясняется релаксационными свойствами материалов. 

Их учет приводит к невозможности мгновенного принятия граничного условия первого рода на стенке (при $\xi =1$) "--- процесс их установления занимает некоторый начальный диапазон времени $0\le \operatorname{\sf Fo}\le

\operatorname{\sf Fo}^\ast $. 

И, в частности, при $\operatorname{\sf Fo}_ 1 =\operatorname{\sf Fo}_2 =0.15$ время принятия граничного условия 1-го рода $\operatorname{\sf Fo}^\ast =0.5$ (см. рис.~\ref{kud:pic24}). 

Отсюда следует, что граничные условия первого рода (тепловой удар) не могут быть установлены мгновенно ни при каких условиях теплообмена с~окружающей средой. 

Следовательно, в~реальных процессах коэффициент теплоотдачи ограничен некоторой постоянной величиной, определяемой

релаксационными свойствами материала.



На рис.~\ref{kud:pic25} приведены результаты расчетов в случае когда $\gamma \ne 0$. Их

анализ позволяет заключить, что при малых значениях $\operatorname{\sf Fo}_1

=\operatorname{\sf Fo}_2 $ величина значений $\gamma $ оказывает незначительное влияние

на время задержки воспламенения. С увеличением значений коэффициентов

релаксации $\operatorname{\sf Fo}_1 $ и $\operatorname{\sf Fo}_2 $ его влияние значительно

возрастает.



%\begin{figure}[h!]\centering \small

%\includegraphics[scale=0.8]{kudinov25.eps}

%\caption{Изменение температуры во времени ($\xi = 0,9$):  1 "--- $\operatorname{\sf Fo}_1 =\operatorname{\sf Fo}_2 =0,0015$; $\gamma =0$; 2 "--- $\operatorname{\sf Fo}_1 =\operatorname{\sf Fo}_2 =0,0015$;

%$\gamma =10^{11}$; 3 "--- $\operatorname{\sf Fo}_1 =\operatorname{\sf Fo}_2 =0,015$; $\gamma =0$; 4 "--- $\operatorname{\sf Fo}_1 =\operatorname{\sf Fo}_2 =0,015$; $\gamma =10^{12}$; 5 "--- $\operatorname{\sf Fo}_1 =\operatorname{\sf Fo}_2 =0,15$; $\gamma =0$; 6 "--- $\operatorname{\sf Fo}_1 =\operatorname{\sf Fo}_2 =0,15$; $\gamma =10^{13}$}

%\label{kud:pic25}

%\end{figure}



\smallskip



\Section{Расчет критических условий теплового взрыва с учетом релаксационных свойств материалов}

Известные математические модели с нелинейными источниками теплоты,

допускающие возникновение условий неограниченного возрастания температуры

(теплового взрыва), основаны на параболических уравнениях, при выводе

которых не учитывается пространственно"=временная нелокальность

(релаксационные свойства материалов). Такие модели в области малых значений

временной и пространственной переменных неадекватно описывают реальные

физические процессы, что связано с локальностью описывающих эти процессы

дифференциальных уравнений, выведенных исходя из принципа локального

термодинамического равновесия и гипотезы сплошной среды [\citen{kud:bib1}--\citen{kud:bib5}]. Подобные

модели допускают бесконечные скорости распространения теплоты и бесконечные

значения тепловых потоков (например, при тепловом ударе на поверхности). В

работах [\citen{kud:bib3}--\citen{kud:bib5}, \citen{kud:bib9}--\citen{kud:bib13}] приводятся принципы построения моделей, свободных

от указанных недостатков, что связано с учетом локальной неравновесности

процессов. В настоящей работе указанные принципы распространены на модели с

нелинейным внутренним источником теплоты.



В качестве конкретного примера рассмотрим задачу теплопроводности для

бесконечной пластины при неоднородных граничных условиях третьего рода с

линейно зависящим от температуры внутренним источником теплоты $\omega

(T)=\omega _0 (1+\beta T)$ в следующей математической постановке:



\begin{gather}

\label{kud:eq46}

(1+\gamma )\frac{\partial \Theta }{\partial \operatorname{\sf Fo}}+\operatorname{\sf Fo}_1 \frac{\partial ^2\Theta }{\partial \operatorname{\sf Fo}^2}=\frac{\partial ^2\Theta }{\partial \xi ^2}+\operatorname{\sf Fo}_2 \frac{\partial ^3\Theta }{\partial \xi^2\partial \operatorname{\sf Fo}}+ \operatorname{\sf Po}_1 \Theta + \operatorname{\sf Po}_2;\\

\label{kud:eq49}

\Theta (\xi, 0)  =0,

\quad

\frac{\partial \Theta (\xi, {\sf Fo}) }{\partial \operatorname{\sf Fo}} \Big|_{\operatorname{\sf Fo}=0} =0;\\

\Bigl[

\frac{\partial \Theta (\xi, {\sf Fo}) }{\partial \xi} +\operatorname{\sf Fo}_2 

\frac{ \partial^2\Theta (\xi, {\sf Fo})}{\partial \xi \partial \operatorname{\sf Fo}}- \hspace{5cm}

\\ \hspace{3cm} -

{\sf Bi}_ 1 \Bigl( \Delta T_1 +

\Theta (\xi, {\sf Fo}) +\operatorname{\sf Fo}_1 \frac{\partial \Theta (\xi, {\sf Fo})} {\partial \operatorname{\sf Fo} } \Bigl)

\Bigr]_{\xi =0} =0;

\label{kud:eq49}

\\

\Bigl[

\frac{\partial \Theta (\xi, {\sf Fo}) }{\partial \xi} +

\operatorname{\sf Fo}_2 \frac{ \partial^2\Theta (\xi, {\sf Fo})} {\partial \xi \partial \operatorname{\sf Fo}}- 

\hspace{5cm}

\\ \hspace{3cm} -

{\sf Bi}_2 \Bigl(

\Delta T_2 -\Theta (\xi, {\sf Fo}) -

\operatorname{\sf Fo}_1 \frac{ \partial \Theta (\xi, {\sf Fo})} {\partial \operatorname{\sf Fo}} \Bigr)\Bigr]_{\xi =1} =0,

\label{kud:eq49}

\end{gather}

где 

$$

\Theta =\frac{T-T_0 }{T_{1} -T_0}, \quad 

\operatorname{\sf Fo}=\frac{at}{\delta ^2}, \quad  

\xi =\frac{x}{\delta } , \quad 

{\sf  Po}_1 =\frac{\omega _0 \delta^2\beta }{ac\rho }, \quad  {\sf Po}_2 = {\sf Po}_1

\frac{1+\beta T_0 }{\beta \Delta T},$$

$$

\Delta T=T_{1} -T_0, \quad 

\operatorname{\sf Fo}_1 =\frac{a\tau _1 }{\delta ^2}, \quad 

\operatorname{\sf Fo}_2 =\frac{a\tau _2 }{\delta ^2}, \quad 

{\sf  Bi}_1 =\frac{\alpha_1 \delta }{\lambda }, \quad 

{\sf Bi}_2 =\frac{\alpha_2 \delta }{\lambda },$$ 

$$\gamma =\frac{\tau _1 \omega _0 \beta }{c\rho }, \quad \Delta T_1 =\frac{T_0 -T_{1}}{\Delta T}, 

\quad  \Delta T_2 =\frac{T_{2} -T_0 }{\Delta T};$$

$T$ "--- температура;

$T_0$ "--- начальная температура; $T_{1}$, $T_{2}$ "---

температуры сред; $\alpha_1$, $\alpha_2$ "--- коэффициенты

теплоотдачи; $\delta $ "--- толщина пластины; $\rho $ "---

плотность; $c$ "--- теплоемкость;  $a$ "--- коэффициент температуропроводности; $\omega_0$ "---

мощность источника теплоты при начальной температуре; $\beta $ "---

коэффициент; $\tau _1$, $\tau_2$ "--- коэффициенты релаксации

соответственно теплового потока и градиента температуры в формуле закона

Фурье; $t$ "--- время; $x$ "--- координата.



Вывод уравнения \eqref{kud:eq46} базируется на формуле закона Фурье \eqref{kud:eq9}, в которой

учитывается ускорение во времени теплового потока и градиента температуры, и

уравнении теплового баланса \eqref{kud:eq45}.



Задача \eqref{kud:eq46}--\eqref{kud:eq49} решалась методом конечных разностей  \cite{kud:mc:bib6}. 

Анализ полученных результатов (при $ {\sf Bi}_1 =5$; $ {\sf Bi}_2 =10$; $\gamma =0$) позволяет

заключить, что при $\beta =0.001$ для любых $\operatorname{\sf Fo}_1$ и

$\operatorname{\sf Fo}_2$ неограниченного возрастания температуры (ситуация теплового

взрыва) не происходит и при $\operatorname{\sf Fo}\ge 0.42$ наблюдается установившееся

(стационарное) состояние, при котором теплоприток от источника теплоты

компенсируется теплоотводом в окружающую среду (см. рис.~\ref{kud:pic26} и кривую {\sl 3} на рис.~\ref{kud:pic27}).



\begin{figure}[h!]



\centering \footnotesize



\begin{tabular}{lp{4.9cm}}

\includegraphics[scale=0.8]{kudinov26}~\,&

\vspace{-5.25cm}

\caption{Распределение температуры $\Theta$ для различных чисел $\sf Fo$, когда $\beta =0.001$, $\operatorname{\sf Fo}_ 1 =\operatorname{\sf Fo}_2 =0.1$

\label{kud:pic26}}

\\

&\vspace{-3.8cm} 

[Figure~\ref{kud:pic26}. The temperature distribution for different Fourier numbers ({\sf Fo})  when \centerline{$\beta =0.001$, ${\operatorname{\sf Fo}_ 1 =\operatorname{\sf Fo}_2 =0.1}$]}

\\



\end{tabular}







\begin{tabular}{lp{4.9cm}}

\includegraphics[scale=0.8]{kudinov27}~\,&

\vspace{-5.25cm}

\caption{Изменение температуры $\Theta$, когда $\xi = 0.9$ и

$\beta =0.1$, $\operatorname{\sf Fo}_ 1 =\operatorname{\sf Fo}_2 = 10^{-3}$ (линия~{\sl 1}\/),  

$\beta \hm =0.1$, $\operatorname{\sf Fo}_ 1 =\operatorname{\sf Fo}_2 =0.05$ (линия~{\sl 2}\/),  

$\beta =0.001$, $\operatorname{\sf Fo}_ 1 =\operatorname{\sf Fo}_2 =0.1$ (линия~{\sl 3}\/)

\label{kud:pic27}}

\\ 

&\vspace{-2.8cm}

[Figure~\ref{kud:pic27}. 

The~change of the temperature ($\Theta$) when $\xi = 0.9$, and 

$\beta =0.1$, $\operatorname{\sf Fo}_ 1 =\operatorname{\sf Fo}_2 = 10^{-3}$ (curve~{\sl 1}\/), and  

$\beta  =0.1$, $\operatorname{\sf Fo}_ 1 =\operatorname{\sf Fo}_2 =0.05$ (curve~{\sl 2}\/),  and

$\beta =0.001$, $\operatorname{\sf Fo}_ 1 =\operatorname{\sf Fo}_2 =0.1$ \centerline{(curve~{\sl 3}\/)]}

\\

\end{tabular}





\begin{tabular}{lp{4.9cm}}

\includegraphics[scale=0.8]{kudinov28}~\,&

\vspace{-5.25cm}

\caption{Изменение температуры $\Theta$, когда $\xi = 0.9$, $\operatorname{\sf Fo}_ 1 =\operatorname{\sf Fo}_2 \hm = 0.1$, $\beta =0.1$ и $\gamma=0$  (линия~{\sl 1}\/),  $\gamma=2.5$ (линия~{\sl 2}\/)

\label{kud:pic28}}

\\ 

&\vspace{-3.8cm}

[Figure~\ref{kud:pic28}. 

The~change of the temperature ($\Theta$) when $\xi = 0.9$, $\operatorname{\sf Fo}_ 1 =\operatorname{\sf Fo}_2  = 0.1$, and 

$\gamma=0$  \centerline{(curve~{\sl 1}\/), and  

$\gamma=2.5$  (curve~{\sl 2}\/)]}

\\

\end{tabular}



\vspace{-10mm}



\end{figure}







С~увеличением $\beta$ (например при $\beta =0.1$) тепловой взрыв

возникает при любых $\operatorname{\sf Fo}_1 $ и $\operatorname{\sf Fo}_2$, причем 

время его задержки существенно зависит от их величины. 

Так, при увеличении $\operatorname{\sf Fo}_1 =\operatorname{\sf Fo}_2$ от их значения, равного $10^{-3}$, до величины $0.05$ время задержки теплового взрыва в~точке $\xi = 0.9$ возрастает

почти в~два раза (см. кривые {\sl 1}, {\sl 2} на рис.~\ref{kud:pic27}). 

Причиной его задержки является тот факт, что при учете релаксационных свойств материала тепловые

потоки на границах в некотором диапазоне времени возрастают от нуля до

значения, определяемого граничными условиями третьего рода. 

В~силу учета релаксационных свойств материала граничные условия (независимо от их вида "---

первого, второго, третьего) могут быть выполнены лишь по истечении некоторого начального диапазона времени $0\le \operatorname{\sf Fo}\le \operatorname{\sf Fo}^{\ast}$, где $\operatorname{\sf Fo}^{\ast}$ "--- время, при котором

начинает выполняться заданное граничное условие краевой задачи.

Следовательно, в определенном диапазоне начального временного участка

прогрев пластины происходит в условиях ограниченного теплообмена с

окружающей средой, участвующей в прогреве конструкции (при $T_{1}>T_0$, $T_{2}>T_0$), что и является сдерживающим

фактором процесса теплового взрыва.



Было выполнено также исследование влияния релаксационных параметров на время

задержки теплового взрыва в случае полного отсутствия теплообмена с~окружающей средой, то есть, когда температура тела определяется действием

лишь внутреннего источника теплоты. Полученные результаты позволяют

заключить, что время задержки теплового взрыва в данном случае существенно

увеличивается, что объясняется учетом ускорения во времени теплового потока

и~градиента температуры.



На рис.~\ref{kud:pic28} приведены результаты расчетов изменения температуры во времени в~точке $\xi =0.9$ при $\gamma =0$ и $\gamma = 2.5$ ($\beta = 0.1$, $\operatorname{\sf Fo}_ 1 =\operatorname{\sf Fo}_2 =0.1$, $ {\sf Bi}_1 =5$, $ {\sf Bi}_2 =10$). 

Их анализ позволяет заключить, что величина $\gamma $ оказывает определенное влияние на время задержки теплового взрыва, которое возрастает с увеличением $\gamma $.



%\begin{figure}[h!]\centering \small

%\includegraphics[scale=0.8]{kudinov28}

%\caption{Изменение температуры при $\xi =0,9$: $\operatorname{\sf Fo}_ 1 =\operatorname{\sf Fo}_2 =0,1$;

%$\beta =0,1$: 1 "--- $\gamma =0$; 2 "--- $\gamma = 2,5$ }

%\label{kud:pic28}

%

%

%

%\end{figure}



\smallskip



\Section{Экспериментальные исследования продольных колебаний \linebreak стержня}

Экспериментальные исследования продольных колебаний стержня были выполнены

на специальном стенде. В процессе испытаний закрепленный на одном из торцов

стержень длиной 30~см растягивался на 0.1~мм. 

После сброса нагрузки с~помощью специальной измерительной системы фиксировались колебания свободного

торца стержня во времени. 

Результаты исследований приведены на рис.~\ref{kud:pic29}. 

Их анализ позволяет заключить, что концевое сечение стержня участвует в~двух

колебательных процессах "--- с~малой амплитудой и~большой частотой (временной

участок $a$ на рис.~\ref{kud:pic29}) и одновременно "--- с~большой амплитудой и~малой

частотой (временной участок $b$ на рис.~\ref{kud:pic29}).



\pagebreak



Полученные экспериментальные данные качественно полностью согласуются с~результатами теоретических исследований по разработанным моделям колебаний упругих тел с~учетом пространственно"=временной нелокальности.



\begin{figure}[t!]\centering \footnotesize

\includegraphics[scale=0.8]{kudinov29}

\caption{Перемещения свободного торца стержня 

\label{kud:pic29}}



[Figure~\ref{kud:pic29}. The displacements of free end of the rod]



\vspace{-6mm}



\end{figure}



Количественное сравнение экспериментальных и теоретических исследований

показывает, что при $\operatorname{\sf Fo}_ 1 =\operatorname{\sf Fo}_2 =10$ расчетная

частота колебаний составляет 2\,380~Гц, а экспериментальная "--- 2\,400~Гц.

Расчетная амплитуда колебаний в диапазоне от 0 до 0.01~с равна

0.0033~мм, средняя экспериментальная амплитуда в этом же диапазоне времени

составляет 0.0028~мм.



\smallskip



\Section[n]{Выводы}  Сформулируем основные выводы по работе.

\begin{itemize}



\item[1.] С учетом пространственно"=временн\'{o}й нелокальности разработана математическая модель теплопереноса, основанная на представлении классической формулы закона Фурье для теплового потока в ряд по степеням малых параметров коэффициентов релаксации теплового потока и~скалярной величины градиента температуры. На основе этой формулы и уравнения теплового баланса получено дифференциальное уравнение локально-неравновесного теплообмена, включающее производные высокого порядка, в том числе и смешанные производные.



\item[2.] Разработаны методы получения точных аналитических решений краевых задач локально"=неравновесного теплообмена, основанные на использовании метода разделения переменных и операционных методов. 

Исследование полученных решений позволило сделать вывод об отсутствии скачков температуры внутри рассматриваемых областей, а также отрицательных значений температур в обратной тепловой волне. Скачки температур, наблюдающиеся лишь в бесконечно малой окрестности граничных точек на начальном участке временн\'{o}й переменной, свидетельствуют о физической невозможности мгновенного задания граничных условий первого рода.



\item[3.] С учетом пространственно"=временн\'{o}й нелокальности разработана математическая модель нестационарного неравновесного теплообмена при ламинарном течении жидкости в цилиндрической трубе. Анализ численного решения задачи (метод прогонки) позволил сделать заключение об отсутствии скачков температуры внутри движущейся жидкости, а~также отрицательных ее значений во всем диапазоне времени нестационарного процесса. Выполненные исследования показали невозможность мгновенного охлаждения жидкости на стенке "--- температура жидкости принимает температуру стенки за некоторый конечный интервал времени.



\item[4.] С учетом пространственно"=временн\'{o}й нелокальности разработана математическая модель затухающих колебаний упругих тел (стержней, пружин, упругих жидкостей и проч.), основанная на представлении классической формулы закона Гука для нормального напряжения по степеням малых параметров коэффициентов релаксации напряжения и~деформации (скалярной величины градиента скорости). На основе этой формулы и уравнения равновесия (уравнение второго закона Ньютона) получено дифференциальное уравнение локально-неравновесных колебаний упругих тел. Исследование его точных аналитических решений позволило заключить, что использование релаксационных составляющих напряжений и деформаций в формуле Гука приводит к устранению скачкообразного изменения напряжений в процессе колебаний среды, выведенной из состояния равновесия.



\item[5.]  Результаты теоретических исследований колебаний закрепленного на одном из торцов стержня сравнивались с данными натурных экспериментов. Стержень подвергался растяжению на 0.1~мм и после сброса нагрузки с помощью лазерного интерферометра определялись его продольные колебания во времени. Результаты экспериментальных исследований подтвердили результаты точного аналитического решения, согласно которому каждая точка стержня участвует в двух колебательных процессах. В одном из них совершаются высокочастотные и низкоамплитудные колебания, а в другом "--- низкочастотные и высокоамплитудные. 

На основе полученного точного аналитического решения и экспериментальных данных  путем решения обратной задачи были найдены коэффициенты релаксации напряжения и градиента перемещения.



\item[6.] Сравнение результатов теоретических исследований с экспериментальными данными было выполнено также и для жидкости (масло) в условиях гидравлического удара. 

С использованием точного аналитического решения и результатов экспериментов, выполненных Е.~Л.~Холмбоу и~В.~Т.~Руло~\cite{kud:bib16}, путем решения обратной задачи были найдены коэффициенты релаксации касательного напряжения и градиента перемещения, которые оказались равными $\tau_1 =0.0108$~c; $\tau _2 =0.0116$~c. Удовлетворительное совпадение результатов теоретических и экспериментальных исследований при найденных таким путем коэффициентах релаксации как для стержня, так и для несжимаемой жидкости  свидетельствует о физической правомерности представления формул законов Гука и Ньютона с~учетом ускорения во времени напряжений (нормального и касательного) и градиентов соответствующих величин (перемещений и скорости).



\item[7.] С учетом пространственно"=временн\'{o}й нелокальности разработана математическая модель динамической термоупругости, основанная на разложении классической формулы закона Гука в ряд по степеням малых параметров коэффициентов релаксации напряжений и деформаций при использовании локально-неравновесной модели теплопереноса. На основе метода разделения переменных  получено точное аналитическое решение краевой задачи локально-неравновесной динамической термоупругости для свободной бесконечной пластины при симметричных граничных условиях первого рода температурной задачи.





\end{itemize}



\bigskip





\AdditionalContent{Конкурирующие интересы}{У меня нет конкурирующих интересов.} 

\AdditionalContent{Авторская ответственность}{Я несу полную ответственность за предоставление окончательной версии рукописи в печать. Окончательная версия рукописи мною одобрена.} 

\AdditionalContent{Финансирование}{Работа поддержана Российским фондом фундаментальных исследований (проект №~18--38--00029\_мол\_а).}

%\AdditionalContent{Благодарность}{Авторы благодарны рецензенту за тщательное прочтение статьи и~ценные предложения и комментарии.}





%\thanks{Исследование выполнено при финансовой поддержке РФФИ в рамках научного

%проекта \No 16-38-00059 мол{\_}а.}



%\thanks{{\bf ORCID}

%

%{\bf Игорь Васильевич Кудинов:}  \url{http://orcid.org/0000-0002-9422-0367}

%

%}



\RusRef

\begin{thebibliography}{99}



\Bibitem{kud:bib1}

\by de Groot~S.~R., Mazur~P. 

\book Non-Equilibrium Thermodynamics

\publaddr Amsterdam

\publ North-Holland

\yr 1962

\totalpages x+510



\Bibitem{kud:bib2}

\serial  Ingenieurwissenschaftliche Bibliothek / Engineering Science Library

\yr 1970

\book Non-equilibrium Thermodynamics. Field Theory and Variational Principles

\by Gyarmati~I.

\crossref{10.1007/978-3-642-51067-0} 

\publ Springer 

\publaddr Berlin

\totalpages xi+184~pp \nofrills





\Bibitem{kud:W.Day}

\serial  Springer Tracts in Natural Philosophy

\vol 22 

\book The Thermodynamics of Simple Materials with Fading Memory

\by Day~W.~A.

\publ Springer 

\publaddr Berlin

\yr 1972

\totalpages x+135~pp \nofrills

\crossref{10.1007/978-3-642-65318-6}



\Bibitem{kud:Brankov}

\by Petrov~N., Brankov~G.

\yr 1982 

\book Modern Problems of Thermodynamics

\publaddr Sofia

\publ Bulgarian Academy of Sciences

\lang In Bulgarian



\Bibitem{kud:Lebon}

\paper Extended irreversible thermodynamics

\by Jou~D.,  Casas-V\'azquez~J., Lebon~G.

\yr 1988 

\jour Rep. Prog. Phys. 

\vol 51 

\issue 8

\pages 1105--1179

\crossref{10.1088/0034-4885/51/8/002}



\Bibitem{kud:bib5}

\by Jou~D.,  Casas-V\'azquez~J., Lebon~G.

\book Extended Irreversible Thermodynamics

\publaddr Netherlands

\publ     Springer 

\yr 2010

\totalpages xviii+483~pp \nofrills

\crossref{10.1007/978-90-481-3074-0}



\RBibitem{kud:bib3}

\by Соболев~С.~Л.

\paper Процессы переноса и бегущие волны в локально"=неравновесных системах

\jour УФН

\vol 161

\issue 3

\yr 1991

\pages 5--29

\crossref{10.3367/UFNr.0161.199103b.0005}



\RBibitem{kud:bib44}

\by Соболев~С.~Л.

\paper Локально"=неравновесные модели процессов переноса

\jour УФН

\vol 167

\issue 10

\yr 1997

\pages 1095--1106

\crossref{10.3367/UFNr.0167.199710f.1095}



\Bibitem{kud:bib42}

\crossref{10.1016/j.ijheatmasstransfer.2016.12.051}

\by Sobolev~S.~L.

\paper Discrete space-time model for heat conduction: Application to size-dependent thermal conductivity in nano-films

\jour Int. J. Heat Mass Transfer

\vol 108

\yr 2017

\issueinfo Part A

\pages 933--939



\Bibitem{kud:bib43}

\by Sobolev~S.~L.

\paper Nonlocal two-temperature model: Application to heat transport in metals irradiated by ultrashort laser pulses

\jour Int. J. Heat Mass Transfer

\vol 94

\yr 2016

\pages 138–144

\crossref{10.1016/j.ijheatmasstransfer.2015.11.075}





\Bibitem{kud:Maxwell}

\paper On the Dynamical Theory of Gases

\by  Maxwell~J.~C. 

\jour Phil. Trans. Royal Soc. London

\yr 1867

\vol 157

\pages  49–88

\moreref 

\crossref{10.1016/b978-0-08-011870-3.50008-8} \nofrills

\moreref 

\crossref{10.1016/b978-0-08-011869-7.50008-4} \nofrills

\moreref 

\crossref{10.1017/cbo9780511710377.003} \nofrills

\moreref 

\crossref{10.1142/9781848161337_0014}. \nofrills









\Bibitem{kud:bib41}

\by Bubnov~V.~A.

\paper More concepts in the theory of heat

\jour Int. J. Heat Mass Transfer

\yr 1976

\vol 19

\issue 2

\pages 175--184

\crossref{10.1016/0017-9310(76)90110-1}



\Bibitem{kud:Mikic}

\crossref{10.1016/0017-9310(67)90059-2}

\paper A model rate equation for transient thermal conduction

\by   Mikic~B.~B.

\jour Int. J. Heat Mass Transfer

\vol 10

\issue 12

\yr 1967

\pages 1899--1904





\Bibitem{kud:Urushev}	

\paper 	Temperature waves and the Boltzmann kinetic equation for phonons

\by 	Urushev~D., Borisov~M., Vavrek~A.

\jour 	Bulgar. J.~Phys.

\vol 15

\issue 6

\yr 1988

\pages 564--575



\Bibitem{kud:Gurtin}	

\yr 1968

\vol 31

\issue 2

\pages 113–126

\paper A general theory of heat conduction with finite wave speeds

\by  Gurtin~M.~E., Pipkin~A.~C.

\crossref{10.1007/BF00281373}

\jour Arch. Rational Mech. Anal. 





\Bibitem{kud:Nunziato}

\by  Nunziato~J.~W.

\paper On heat conduction in materials with memory 

\jour Quart. Appl. Math. 

\vol 29 

\issue 2

\yr  1971

\pages 187--204

\crossref{10.1090/qam/295683}





\Bibitem{kud:Weymann}

\paper Finite Speed of Propagation in Heat Conduction, Diffusion, and Viscous Shear Motion

\by Weymann~H.~D. 

\jour Am.  J.  Phys.

\vol 35

\issue 6

\pages 488--496

\yr 1967

\crossref{10.1119/1.1974155}



\Bibitem{kud:Taitel}

\jour Int. J. Heat Mass Transfer

\vol 15

\issue 2

\yr 1972

\pages 369--371

\paper  On the parabolic, hyperbolic and discrete formulation of the heat conduction equation

\by Taitel~Y. 

\crossref{10.1016/0017-9310(72)90085-3}



\RBibitem{kud:bib31}

\by Леванов~Е.~И., Сотский~Е.~Н.

\paper Теплоперенос с учетом релаксации теплового потока

\inbook Математическое моделирование (нелинейные дифференциальные уравнения математической физики)

\publaddr М.

\publ Наука

\yr 1987

\pages 155--196



\RBibitem{kud:bib32}

\by Черешнев~С.~Л., Генич~А.~П., Куликов~С.~В., Манелис~Г.~Б.

\preprint Эффекты поступательной неравновесности в~ударных волнах в~газах

\preprintinfo Препринт; АН СССР, Отд-ние Ин-та хим. физики

\publaddr Черноголовка

\publ ОИХФ

\yr 1988

\totalpages 71



\Bibitem{kud:MacDonald}

\by MacDonald~R.~A., Tsai~D.~H.

\jour Phys. Rep.

\vol 46

\issue 1

\yr 1978

\pages 1--41

\paper Molecular dynamical calculations of energy transport in crystalline solids

\crossref{10.1016/0370-1573(78)90006-6}





\RBibitem{kud:bib6}

\by Лыков~А.~В.

\book Теория теплопроводности

\publaddr М.

\publ Высшая школа

\yr 1967

\totalpages 600



\RBibitem{kud:bib7}

\by Лыков~А.~В.

\book Тепломассообмен

\publaddr М.

\publ Энергия

\yr 1978

\totalpages 480



\RBibitem{kud:bib8}

\by Лыков~А.~В.

\paper Применение методов термодинамики необратимых процессов к исследованию тепло- и массообмена

\jour Инженерно"=физический журнал

\vol 9

\issue 3

\yr 1965

\pages 287--304



\RBibitem{kud:bib14}

\by Филин~А.~П.

\book Прикладная механика твердого деформируемого тела

\vol 1

\publaddr М.

\publ Наука

\yr 1975

\totalpages 832



\RBibitem{kud:bib9}

\by Кудинов~И.~В., Кудинов~В.~А.

\paper Математическая модель локально-неравновесного теплопереноса с учетом пространственно-временной нелокальности

\jour Инженерно"=физический журнал

\yr 2015

\vol 88

\issue 2

\pages 393--408

\elib{https://elibrary.ru/item.asp?id=23274506}



\RBibitem{kud:bib10}

\by Кудинов~И.~В., Кудинов~В.~А.

\book Аналитические решения параболических и гиперболических уравнений тепломассопереноса

\publaddr М.

\publ Инфра-М

\yr 2013

\totalpages 391



\RBibitem{kud:bib11}

\by Кудинов~В.~А., Кудинов~И.~В.

\paper Получение и анализ точного аналитического решения гиперболического уравнения теплопроводности для плоской стенки

\jour ТВТ

\yr 2012

\vol 50

\issue 1

\pages 118--125

\mathnet{http://mi.mathnet.ru/tvt237}

\elib{http://elibrary.ru/item.asp?id=17313267}

%\transl

%\jour High Temperature

%\yr 2012

%\vol 50

%\issue 1

%\pages 112--119

%\crossref{10.1134/S0018151X12010105}

%\isi{http://gateway.isiknowledge.com/gateway/Gateway.cgi?GWVersion=2&SrcApp=PARTNER_APP&SrcAuth=LinksAMR&DestLinkType=FullRecord&DestApp=ALL_WOS&KeyUT=000303513700016}

%\elib{http://elibrary.ru/item.asp?id=17986087}

%\scopus{http://www.scopus.com/record/display.url?origin=inward&eid=2-s2.0-84861377979}



\RBibitem{kud:bib12}

\by Кудинов~В.~А., Кудинов~И.~В.

\paper Исследование теплопроводности с учетом конечной скорости распространения теплоты

\jour ТВТ

\yr 2013

\vol 51

\issue 2

\pages 301--310

\mathnet{http://mi.mathnet.ru/tvt85}

\elib{http://elibrary.ru/item.asp?id=18822194}

%\transl

%\jour High Temperature

%\yr 2013

%\vol 51

%\issue 2

%\pages 268--276

%\crossref{10.1134/S0018151X1204013X}

%\isi{http://gateway.isiknowledge.com/gateway/Gateway.cgi?GWVersion=2&SrcApp=PARTNER_APP&SrcAuth=LinksAMR&DestLinkType=FullRecord&DestApp=ALL_WOS&KeyUT=000320157900017}

%\elib{http://elibrary.ru/item.asp?id=20445720}

%\scopus{http://www.scopus.com/record/display.url?origin=inward&eid=2-s2.0-84880671446}





\RBibitem{kud:bib13}

\by Кудинов~В.~А., Кудинов~И.~В.

\paper Получение точных аналитических решений гиперболических уравнений движения при разгонном течении Куэтта

\jour  Изв. РАН. Энергетика

\issue 1

\yr 2012

\pages 119--133







\RBibitem{kud:bib15}

\by Чарный~И.~А.

\book Неустановившееся движение реальной жидкости в трубах

\publaddr М.

\publ Недра

\yr 1975

\totalpages 296



\Bibitem{kud:bib16}

\crossref{10.1115/1.3609549}

\paper The effect of viscous shear on transients in liquid lines

\by 	Holmboe~E.~L., Rouleau~W.~T.

\yr 	1967

\jour J.~Basic Eng.

\vol 89

\issue 1

\pages 174--180 



\RBibitem{kud:bib17}

\by Еремин~А.~В., Кудинов~И.~В., Кудинов~В.~А.

\paper Математическая модель теплообмена в жидкости с учетом ее релаксационных свойств

\jour Изв. РАН. МЖГ

\yr 2016

\issue 1

\pages 33--44

\elib{https://elibrary.ru/item.asp?id=25809209}





\RBibitem{kud:bib18}

\by Даниловская~В.~И.

\paper Динамические температурные напряжения в бесконечной плите

\jour Инженерный журнал

\yr 1961

\vol 1

\issue 4

\pages 86--94



\RBibitem{kud:bib19}

\by Даниловская~В.~И., Зубчанинова~В.~Н.

\paper Температурные поля и напряжения, возникающие в пластинке вследствие потока лучистой энергии

\jour Прикладная механика

\yr 1968

\vol 4

\issue 1

\pages 103--110

%http://www.mathnet.ru/links/997dc0dd4bc289da315e9f4c1dadeda9/tvt_6698_refs_rus.htm

%https://link.springer.com/content/pdf/10.1007/BF00920328.pdf



\Bibitem{kud:bib20}

\book Theory of Thermal Stresses

\by Boley~B.~A., Weiner~J.~H. 

\publaddr New York 

\publ John Wiley

\yr 1960

\totalpages 586



\RBibitem{kud:bib21}

\by Семерак~Ф.~В. Борисенко~О.~И.

\paper Динамическая задача термоупругости для бесконечной пластинки

\jour  Математичні методи та фізико-механічні поля

\yr 1977

\vol 6

\pages 61--63



\RBibitem{kud:bib22}

\by Франк--Каменецкий~Д.~А.

\book  Основы макрокинетики. Диффузия и теплопередача в химической кинетике

\publaddr М.

\publ  Наука

\yr 1967

\totalpages 492



\RBibitem{kud:bib23}

\by Мержанов~А.~Г., Дубовицкий~Ф.~И.

\paper Современное состояние теории теплового взрыва

\jour Усп. хим.

\yr 1966

\vol 35

\issue 4

\pages 656--683

\mathnet{rcr1978}





\RBibitem{kud:mc:bib6}

\by Годунов~С.~К., Рябенький~В.~С.

\book Разностные схемы. Введение в теорию

\publaddr  М.

\publ Наука

\yr 1977

\totalpages 439







\end{thebibliography}



\newpage







\makeentitle



\renewcommand{\baselinestretch}{0.83}



\newpage



\AdditionalContent{Competing interests}{I have no competing interests.} 

\AdditionalContent{Author's Responsibilities}{I take full responsibility for submitting the final manuscript in print. I approved the final version of the manuscript.} 

\AdditionalContent{Funding}{This work was supported by the Russian Foundation for Basic Research (project no. 18--38--00029\_mol\_a).}



\vspace{-2mm}



\begin{thebibliography}{99}



\Bibitem{kud:bib1:em}

\by de Groot~S.~R., Mazur~P. 

\book Non-Equilibrium Thermodynamics

\publaddr Amsterdam

\publ North-Holland

\yr 1962

\totalpages x+510



\Bibitem{kud:bib2:en}

\serial  Ingenieurwissenschaftliche Bibliothek / Engineering Science Library

\yr 1970

\book Non-equilibrium Thermodynamics. Field Theory and Variational Principles

\by Gyarmati~I.

\crossref{10.1007/978-3-642-51067-0} 

\publ Springer 

\publaddr Berlin

\totalpages xi+184~pp \nofrills





\Bibitem{kud:W.Day:en}

\serial  Springer Tracts in Natural Philosophy

\vol 22 

\book The Thermodynamics of Simple Materials with Fading Memory

\by Day~W.~A.

\publ Springer 

\publaddr Berlin

\yr 1972

\totalpages x+135~pp \nofrills

\crossref{10.1007/978-3-642-65318-6}



\Bibitem{kud:Brankov:en}

\by Petrov N., Brankov~G.

\yr 1982 

\book Modern Problems of Thermodynamics

\publaddr Sofia

\publ Bulgarian Academy of Sciences

\lang In Bulgarian



\Bibitem{kud:Lebon:en}

\paper Extended irreversible thermodynamics

\by Jou~D.,  Casas-V\'azquez~J., Lebon~G.

\yr 1988 

\jour Rep. Prog. Phys. 

\vol 51 

\issue 8

\pages 1105--1179

\crossref{10.1088/0034-4885/51/8/002}



\Bibitem{kud:bib5:en}

\by Jou~D.,  Casas-V\'azquez~J., Lebon~G.

\book Extended Irreversible Thermodynamics

\publaddr Netherlands

\publ     Springer 

\yr 2010

\totalpages xviii+483~pp \nofrills

\crossref{10.1007/978-90-481-3074-0}



\Bibitem{kud:bib3:en}

\by Sobolev~S.~L.

\paper Transport processes and traveling waves in systems with local nonequilibrium

\jour Sov. Phys. Usp. 

\vol 34 

\issue 3

\pages 217–229 

\yr 1991

\crossref{10.1070/PU1991v034n03ABEH002348}



\Bibitem{kud:bib44:en}

\by Sobolev~S.~L.

\paper Local non-equilibrium transport models

\jour Phys. Usp. 

\vol 40 

\issue 10

\pages 1043–1053 

\yr 1997

\crossref{10.1070/PU1997v040n10ABEH000292}





\Bibitem{kud:bib42:en}

\crossref{10.1016/j.ijheatmasstransfer.2016.12.051}

\by Sobolev~S.~L.

\paper Discrete space-time model for heat conduction: Application to size-dependent thermal conductivity in nano-films

\jour Int. J. Heat Mass Transfer

\vol 108

\yr 2017

\issueinfo Part A

\pages 933--939



\Bibitem{kud:bib43:en}

\by Sobolev~S.~L.

\paper Nonlocal two-temperature model: Application to heat transport in metals irradiated by ultrashort laser pulses

\jour Int. J.~Heat Mass Transfer

\vol 94

\yr 2016

\pages 138–144

\crossref{10.1016/j.ijheatmasstransfer.2015.11.075}





\Bibitem{kud:Maxwell:en}

\paper On the Dynamical Theory of Gases

\by  Maxwell~J.~C. 

\jour Phil. Trans. Royal Soc. London

\yr 1867

\vol 157

\pages  49–88

\moreref 

\crossref{10.1016/b978-0-08-011870-3.50008-8} \nofrills

\moreref 

\crossref{10.1016/b978-0-08-011869-7.50008-4} \nofrills

\moreref 

\crossref{10.1017/cbo9780511710377.003} \nofrills

\moreref 

\crossref{10.1142/9781848161337_0014}. \nofrills









\Bibitem{kud:bib41:en}

\by Bubnov~V.~A.

\paper More concepts in the theory of heat

\jour Int. J.~Heat Mass Transfer

\yr 1976

\vol 19

\issue 2

\pages 175--184

\crossref{10.1016/0017-9310(76)90110-1}



\Bibitem{kud:Mikic:en}

\crossref{10.1016/0017-9310(67)90059-2}

\paper A model rate equation for transient thermal conduction

\by   Mikic~B.~B.

\jour Int. J.~Heat Mass Transfer

\vol 10

\issue 12

\yr 1967

\pages 1899--1904





\Bibitem{kud:Urushev:en}	

\paper 	Temperature waves and the Boltzmann kinetic equation for phonons

\by 	Urushev~D., Borisov~M., Vavrek~A.

\jour 	Bulgar. J.~Phys.

\vol 15

\issue 6

\yr 1988

\pages 564--575



\Bibitem{kud:Gurtin:en}	

\yr 1968

\vol 31

\issue 2

\pages 113–126

\paper A general theory of heat conduction with finite wave speeds

\by  Gurtin~M.~E., Pipkin~A.~C.

\crossref{10.1007/BF00281373}

\jour Arch. Rational Mech. Anal. 





\Bibitem{kud:Nunziato:en}

\by  Nunziato~J.~W.

\paper On heat conduction in materials with memory 

\jour Quart. Appl. Math. 

\vol 29 

\issue 2

\yr  1971

\pages 187--204

\crossref{10.1090/qam/295683}





\Bibitem{kud:Weymann:en}

\paper Finite Speed of Propagation in Heat Conduction, Diffusion, and Viscous Shear Motion

\by Weymann~H.~D. 

\jour Am.  J.~Phys.

\vol 35

\issue 6

\pages 488--496

\yr 1967

\crossref{10.1119/1.1974155}



\Bibitem{kud:Taitel:en}

\jour Int. J. Heat Mass Transfer

\vol 15

\issue 2

\yr 1972

\pages 369--371

\paper  On the parabolic, hyperbolic and discrete formulation of the heat conduction equation

\by Taitel~Y. 

\crossref{10.1016/0017-9310(72)90085-3}



\Bibitem{kud:bib31:en}

\lang In Russian

\by Levanov~E.~I., Sotskiy~E.~N.

\paper Heat transfer with taking into account heat-flux relaxation

\inbook Mathematical Simulation. Nonlinear Differential Equations of Mathematical Physics

\publaddr Moscow

\publ Nauka

\yr 1987

\pages 155--196



\Bibitem{kud:bib32:en}

\lang In Russian

\by Chereshnev~S.~L., Genich~A.~P., Kulikov~S.~V., Manelis~G.~B.

\preprint  Effects of Translational Nonequilibrium in Shock Waves in Gases

\preprintinfo Preprint

\publaddr Chernogolovka

\publ Division of the Institute of Chemical Physics, Academy of Sciences of the USSR

\yr 1988

\totalpages 71







\Bibitem{kud:MacDonald:en}

\by MacDonald~R.~A., Tsai~D.~H.

\jour Phys. Rep.

\vol 46

\issue 1

\yr 1978

\pages 1--41

\paper Molecular dynamical calculations of energy transport in crystalline solids

\crossref{10.1016/0370-1573(78)90006-6}





\Bibitem{kud:bib6:en}

\lang In Russian

\by  Lykov~A.~V.

\book Teoriia teploprovodnosti \rm [Theory of Heat Conduction]

\publaddr Moscow

\publ Vysshaia Shkola

\yr 1967

\totalpages 600



\Bibitem{kud:bib7:en}

\lang In Russian

\by  Lykov~A.~V.

\book Teplomassoobmen \rm [Heat and Mass Transfer]

\publaddr Moscow

\publ Energiia 

\yr 1978

\totalpages 480



\Bibitem{kud:bib8:en}

\jour J.~Eng. Phys. 

\yr 1965

\vol 9

\issue 3

\pages 189–202 

\paper Application of the methods of thermodynamics of irreversible processes to the investigation of heat and mass transfer

\by  Luikov~A.~V.

\crossref{10.1007/BF00828333}





\Bibitem{kud:bib14:en}

\lang In Russian

\by  Filin~A.~P.

\book Prikladnaia mehanika tverdogo deformiruemogo tela \rm [Applied Mechanics of Solid Body]

\vol 1

\publaddr Moscow

\publ Nauka

\yr 1975

\totalpages 832



\Bibitem{kud:bib9:en}

\crossref{10.1007/s10891-015-1206-6}

\by  Kudinov~I.~V.,  Kudinov~V.~A.

\paper Mathematical simulation of the locally nonequilibrium heat transfer in a body with account for its nonlocality in space and time

\jour J.~Eng. Phys. Thermophys.

\vol 88

\issue 2

\pages 406–422 

\yr 2015





\Bibitem{kud:bib10:en}

\lang In Russian

\by Kudinov~I.~V., Kudinov~V.~A.

\book Analiticheskie resheniia parabolicheskikh i giperbolicheskikh uravnenii teplomassoperenosa \rm [Analytical Solutions to Parabolic and Hyperbolic Equations of Heat and Mass Transfer]

\publ Infra-M

\publaddr Moscow

\yr 2013

\totalpages 391



\Bibitem{kud:bib11:en}

\paper One method of reception of the exact analytical decision of the hyperbolic equation of heat conductivity on the basis of use of orthogonal methods

\by  Kudinov~V.~A., Kudinov~I.~V. 

\jour High Temperature

\yr 2012

\vol 50

\issue 1

\pages 112--119

\crossref{10.1134/S0018151X12010105}

\isi{http://gateway.isiknowledge.com/gateway/Gateway.cgi?GWVersion=2&SrcApp=PARTNER_APP&SrcAuth=LinksAMR&DestLinkType=FullRecord&DestApp=ALL_WOS&KeyUT=000303513700016}

\elib{http://elibrary.ru/item.asp?id=17986087}

\scopus{http://www.scopus.com/record/display.url?origin=inward&eid=2-s2.0-84861377979}



\Bibitem{kud:bib12:en}

\paper Studying heat conduction taking into account the finite rate of heat propagation

\by  Kudinov~V.~A.,  Kudinov~I.~V.

\jour High Temperature

\yr 2013

\vol 51

\issue 2

\pages 268--276

\crossref{10.1134/S0018151X1204013X}

\isi{http://gateway.isiknowledge.com/gateway/Gateway.cgi?GWVersion=2&SrcApp=PARTNER_APP&SrcAuth=LinksAMR&DestLinkType=FullRecord&DestApp=ALL_WOS&KeyUT=000320157900017}

\elib{http://elibrary.ru/item.asp?id=20445720}

\scopus{http://www.scopus.com/record/display.url?origin=inward&eid=2-s2.0-84880671446}





\Bibitem{kud:bib13:en}

\lang In Russian

\by Kudinov~V.~A.,  Kudinov~I.~V.

\paper Calculation of Exact Analytic Solutions of Hyperbolic Equations of Motion in the Accelerated Couette Flow

\jour Izv. Ross. Akad. Nauk. Energetika

\issue 1

\yr 2012

\pages 119--133







\Bibitem{kud:bib15:en}

\lang In Russian

\by Charnyi~I.~A.

\book Neustanovivsheesia dvizhenie real'noi zhidkosti v trubakh \rm  [Nonsteady Motion of Real Liquid in Pipes]

\publaddr Moscow

\publ Nedra

\yr 1975

\totalpages 296



\Bibitem{kud:bib16:en}

\crossref{10.1115/1.3609549}

\paper The effect of viscous shear on transients in liquid lines

\by 	Holmboe~E.~L., Rouleau~W.~T.

\yr 	1967

\jour J.~Basic Eng.

\vol 89

\issue 1

\pages 174--180 





\Bibitem{kud:bib17:en}

\paper Mathematical model of heat transfer in a fluid with account for its relaxation properties

\by Eremin~A.~V., Kudinov~V.~A., Kudinov~I.~V.

\jour Fluid Dynamics

\yr 2016

\vol 51

\issue 1

\pages 33--44

\elib{https://elibrary.ru/item.asp?id=27138267}

\crossref{10.1134/S0015462816010051}





\Bibitem{kud:bib18:en}

\lang In Russian

\by  Danilovskaya~V.~I.

\paper Dynamic temperature stresses in an infinite plate

\jour Inzh. Zh.

\yr 1961

\vol 1

\issue 4

\pages 86--94



\Bibitem{kud:bib19:en}

\by Danilovskaya~V.~I., Zubchaninova~V.~N.

\jour Soviet Applied Mechanics 

\yr 1968

\vol 4

\issue 1

\pages 63–66 

\paper Temperature fields and stresses created in a plate by radiant energy

\crossref{10.1007/BF00920328}





\Bibitem{kud:bib20:en}

\book Theory of Thermal Stresses

\by Boley~B.~A., Weiner~J.~H. 

\publaddr New York 

\publ John Wiley

\yr 1960

\totalpages 586



\Bibitem{kud:bib21:en}

\lang In Russian

\by Semerak~F.~V. Borisenko~O.~I.

\paper The dynamic problem of thermoelasticity for an infinite plate

\jour  Matematichni metodi ta fiziko-mekhanichni polia

\yr 1977

\vol 6

\pages 61--63





\Bibitem{kud:bib22:en}

\lang In Russian

\by Frank--Kamenetskii~D.~A.

\book  Osnovy makrokinetiki. Diffuziia i teploperedacha v khimicheskoi kinetike \rm [Basics of macrokinetics. Diffusion and heat transfer in chemical kinetics]

\publaddr Moscow

\publ  Nauka

\yr 1967

\totalpages 492







\Bibitem{kud:bib23:en}

\crossref{10.1070/RC1966v035n04ABEH001462}

\by Merzhanov~A.~G., Dubovitskii~F.~I.

\paper Present state of the theory of thermal explosions

\jour Russ. Chem. Rev.

\yr 1966

\vol 35 

\issue 4

\pages 278–292 





\Bibitem{kud:mc:bib6:en}

\book Theory of Difference Schemes: An Introduction

\by Godunov~S.~K., Ryabenki~V.~S.

\publ Wiley

\publaddr New York

\yr 1964

\totalpages xii+289 







\end{thebibliography}







\renewcommand{\baselinestretch}{0.9}



\newpage



%\end{document}

